\documentclass[11pt,letterpaper]{article}
\usepackage[T1]{fontenc}
\usepackage[utf8]{inputenc}
\usepackage{lmodern,microtype,amsmath,amssymb,amsthm,mathtools}
\usepackage[margin=1.02in,top=.91in,bottom=.96in,headheight=15pt]{geometry}
\usepackage{booktabs,tabularx,array,xcolor,enumitem,fancyhdr,xurl}
\definecolor{ink}{HTML}{173546}
\definecolor{accent}{HTML}{17695E}
\usepackage[colorlinks=true,linkcolor=ink,urlcolor=accent,citecolor=accent]{hyperref}
\hypersetup{pdftitle={Ordinal numbers in intuitionistic primitive recursive arithmetic},pdfsubject={Constructive order types, arithmetic and cofinality}}
\setlength{\emergencystretch}{3.3em}
\setlength{\parskip}{2.5pt plus 1pt}
\setlist{topsep=4pt,itemsep=2pt}
\setcounter{tocdepth}{1}
\pagestyle{fancy}\fancyhf{}\fancyhead[L]{\small\sffamily\color{ink}Ordinal numbers in iPRA}\fancyhead[R]{\small\sffamily\color{ink}Constructive arithmetic}\fancyfoot[C]{\small\thepage}
\renewcommand{\headrulewidth}{.25pt}
\newtheoremstyle{main}{9pt}{4pt}{\itshape}{}{\bfseries\color{ink}}{.}{.5em}{}
\theoremstyle{main}
\newtheorem{theorem}{Theorem}[section]
\newtheorem{lemma}[theorem]{Lemma}
\newtheorem{proposition}[theorem]{Proposition}
\newtheorem{corollary}[theorem]{Corollary}
\theoremstyle{definition}
\newtheorem{example}[theorem]{Example}
\newtheorem{definition}[theorem]{Definition}
\newtheorem{remark}[theorem]{Remark}
\begin{document}
\begin{center}
{\sffamily\small\bfseries\color{accent} CONSTRUCTIVE MATHEMATICS}\par\vspace{11pt}
{\LARGE\bfseries\color{ink} Ordinal numbers in intuitionistic\\[3pt] primitive recursive arithmetic}\par\vspace{9pt}
{\large Order types, Cantor normal forms, and effective cofinality}\par\vspace{8pt}
{\small Research report for Vladimir Reshetnikov\quad $\cdot$\quad November 22, 2026}
\end{center}
\begin{abstract}
A finite linear order determines its order type by counting predecessors. The same first-order theory can reason about an infinite order such as $\omega^\omega$ when its elements have finite canonical names. We construct arithmetic, ordinal division with unique remainder, order isomorphisms for sums, products and finite-support powers, and cofinal sequences for all hereditary Cantor normal forms below $\varepsilon_0$. Given $y<\lambda$ with $\lambda$ a limit, we compute the first index at which the fundamental sequence passes $y$. It is bounded by the largest coefficient anywhere in the hereditary form of $y$. An extension by epsilon constants presents the order below $\varepsilon_\omega$ and identifies its principal-power fixed points. Further applications identify topological derivatives, natural ordinal arithmetic, and the decrease at each hereditary base-change step. The proofs use finite recursion, decidable tests and quantifier-free induction. The formalization section distinguishes the Lean theorems checked here from the remaining work of producing formal iPRA derivations.
\end{abstract}
\tableofcontents
\newpage
\section{What an ordinal statement means in iPRA}
\label{sec:foundation}
The article \emph{Ordinal number} treats ordinals both as transitive sets and as order types. Two examples suggest how to read it arithmetically. The assertion $3=\{0,1,2\}$ can be a literal calculation with four integers coding finite sets. The assertion $\omega^2=\sup_n\omega n$ has a different reading. Every point of $\omega^2$ lies in a finite-numbered block of type $\omega$. Given a point, we can compute the number of a later block. We work out both readings, including the maps and witnesses which distinguish them. The final section gives the source correspondence.

\subsection{Induction and finite data}
Our iPRA is first-order intuitionistic logic with equality, the zero and successor axioms, a symbol and defining equations for each primitive recursive (PR) function, and induction for quantifier-free formulas with numerical parameters. A PR definition is built by composition and ordinary numerical primitive recursion. We use $\mathbb N$ to include zero. Searches below a supplied bound return a first index, or a default outside the range. A predicate whose characteristic function is PR is decidable here. Trichotomy for a coded comparison permits a case split even when the conclusion of the argument contains quantifiers.

Write $\mathrm{All}(0,p)=1$ and $\mathrm{All}(s+1,p)=\mathrm{All}(s,p)\cdot p(s)$ for a fixed Boolean PR test $p$ with arbitrary fixed numerical parameters. A least-failure register in the same loop proves that the output is $1$ exactly when every index below the bound passes. If the output is $0$, the register supplies a failing index. The analogous disjunction scan supplies a first witness. These claims use the loop equations and induction on numerical equalities or bounded tests. For example, the saved failure flag implies $j<s$ and $p(j)=0$, while the forward direction for a selected $i<s$ follows by carrying that fixed index through the loop. Nested bounded quantifiers are consequently shorthand for a single PR test. An induction below using a bounded universal uses that test as induction formula. We introduce universal input variables after proving the instance with them as parameters. We use ordinary intuitionistic quantifiers for conclusions, for hypotheses available at particular arguments, and for schemata over fixed PR definitions. There is no universal PR evaluator for all PR programs.

Here is one consistent realization of the data structures. Code an empty list by $0$ and a cons cell with head $u$ and tail $v$ by
\[
 C(u,v)=1+(u+v)^2+u .
\]
The two components are smaller than this positive code. To decode $c>0$, set $s=\lfloor\sqrt{c-1}\rfloor$ by bounded search, $u=(c-1)-s^2$. Require $u\le s$, and set $v=s-u$. In the permitted range $s^2+u<(s+1)^2$, so this is inverse to construction. Test the pair again on decoding, and use fixed defaults on malformed data. Valid binary trees have valid subtrees in both positions. A list with arbitrary numerical entries needs only its tail to be a valid list. Scans, length, retrieval, concatenation and insertion sort follow by recursion with at most the input code (or the explicit length) many steps. To tabulate $f(t)$ with calls on smaller numerical codes, retain the list $[f(0),\ldots,f(t-1)]$ and retrieve each requested entry by a bounded scan. The history itself is built by one PR recursion. Its potentially large numerical code is allowed. This also implements a fixed finite product of recursions by packing the registers. The same construction works for pairs and triples of subtrees by loops over their decreasing total code or by tables over a fixed box of input codes.

\subsection{The notation and its order}
A term is a finite tree displayed as the list of its children. Compare two lists lexicographically using this same comparison on entries, with the proper prefix smaller. Let $\mathcal O$ consist of the trees for which every list, including every descendant list, is weakly decreasing. Empty is written $0$. Write $[a]$ for a singleton, $1=[0]$, $\omega=[1]$. An expression such as $[2,2,0]$ denotes $\omega^2\cdot 2+1$. Repetitions carry coefficients. The exponent of the largest power in a nonempty term is its first child, also called its degree. A code for a sequence of terms and the list of children of a single term can have the same numerical representation. The indicated domain test specifies which is meant.

\begin{theorem}[Comparison, normal forms, and finite extrema]\label{thm:order}
Validity, equality and comparison are PR. Comparison gives a strict linear order with least element $0$. Each term is uniquely specified by the strictly descending list of its distinct exponents and their positive multiplicities. Every nonempty coded finite family of terms has a witnessed least and greatest member. Its least upper bound is its greatest member. The empty family has least upper bound $0$.
\end{theorem}
\begin{proof}
Compare the empty and nonempty cases directly. For two nonempty cells compare the heads and, only if equal, the tails. The recursive pairs have smaller total code and remain inside the input box. Thus a PR table computes the comparison, including the equality output. A second table tests validity using the two subtrees and the order of adjacent heads. Equality output forces equality of both heads and tails by the same decreasing computation, whence equality of codes. The converse is immediate by the defining clauses. Interchanging inputs reverses the output. For transitivity of $x<y<z$ write their respective heads as $a,b,c$ when all are nonempty. The two comparisons give $a\le b\le c$. If either is strict, transitivity on smaller heads gives $a<c$ (or substitution does so in an equality case). If both are equal, the two strict tail comparisons give the result on smaller tails. The empty cases use the prefix rule. Formally fix a numerical bound $B$ on the three original inputs. Induct on $s$ using the Boolean conjunction of the comparison implications for all triples with components $\le B$ and total code $\le s$. Every subcomparison used in the proof is covered at the preceding step. This is quantifier-free induction, independent of the putative well-ordering of the terms.

Grouping equal adjacent exponents is inverted by repeating each exponent its recorded number of times. A first differing exponent or multiplicity decides comparison, so distinct grouped normal forms remain distinct. A scan of a finite family updates the index of its current extremum. The invariant says its stored value is an input and bounds all scanned entries. For a least upper bound, every upper bound bounds the selected greatest entry. For the empty list every term is an upper bound and $0$ is least.
\end{proof}

\paragraph{Internal and external readings.}
Internally $x<a$ is a decidable numerical predicate on valid codes. We write $I_a$ for this predicate, the strict initial segment of $a$. An isomorphism of displayed presentations means explicit mutually inverse maps preserving the displayed comparisons. A supremum statement may be a universal property quantified over codes. Its proof below usually computes a witness against every smaller candidate. Externally the conventional interpretation of a tree is $|[a_1,\ldots,a_k]|=\omega^{|a_1|}+\cdots+\omega^{|a_k|}$. Classical Cantor normal form identifies precisely the ordinals below $\varepsilon_0$ with these trees. To recall the range argument, the finite heights of trees bound their values by finite towers of $\omega$, and every ordinal below such a tower has a finite CNF whose exponents are below the previous tower. Induction on the tower height represents it. Classical uniqueness and the usual lexicographic comparison give injectivity. This paragraph identifies the classical referents. The arithmetical proofs below work on codes directly. All unqualified term variables from here on satisfy the validity test.

\section{Finite order types and the von Neumann construction}
\label{sec:finite-models}

A finite ordinal has two concrete representatives here. Its order
presentation is the interval $I_n=\{i:i<n\}$ in the usual order, and its
set presentation is a number whose binary digits encode the preceding
ordinals. The passage from a finite relation to the second representative
also gives a finite form of the Mostowski collapse theorem.

\subsection{Classification by rank}

A presentation on $I_n$ consists of a Boolean table $R$ with $n^2$ entries.
It is a strict linear order if the diagonal is false, the relation is
transitive, and exactly one of $R(x,y),R(y,x)$ holds for distinct inputs.
All the tests range over $I_n$. A subset is a mask on $I_n$, and a map is
a list of values indexed by its domain. A presentation with holes in its
domain can first enumerate its occupied positions in numerical order.

\paragraph{Least elements and sorted enumeration.}
For each nonempty mask $D$, scan the points of $I_n$ in numerical order.
Keep a flag and the least point of $D$ seen so far, replacing the incumbent
exactly when the new point precedes it in $R$. After scanning $j$ positions,
the flag says whether an occupied position has occurred, and a flagged
incumbent belongs to the scanned portion and precedes or equals every
occupied point there. The empty prefix initializes the invariant. A
replacement preserves it by transitivity, and retaining the incumbent
preserves it by totality. Thus the algorithm returns a least point.
Two least points must coincide by strictness. Reversing the comparison
gives a greatest point.

Start with the full mask. Repeat least-point selection and deletion $n$
times, writing the selected points as $a_0,\ldots,a_{n-1}$. After $j$
successful selections the remaining mask has $n-j$ bits set, the entries
selected so far are increasing, and every selected entry precedes every
remaining entry. Deleting a single selected bit decreases the count by
one, as follows by splitting its defining sum at that index.
Least-point selection preserves the two comparison assertions. Consequently
the procedure succeeds at every prescribed stage, exhausts the domain,
and produces each point exactly once.

Define $r_A(x)=\sum_{y<n}[R(y,x)]$, where square brackets denote a
Boolean value treated as a number. The predecessors of $a_i$ are exactly
$a_0,\ldots,a_{i-1}$. Reordering a Boolean sum by this enumeration
preserves the sum, since deletion subtracts exactly the deleted summand.
Hence $r_A(a_i)=i$. The maps $i\mapsto a_i$ and $r_A$ are inverse
order-isomorphisms. In particular every finite linear order is a
well-order for coded subsets, its first unused ordinal label is $n$, and
its order type equals its cardinality.

\paragraph{Segments, comparison and rigidity.}
If $D$ is downward closed in $A$, let $k$ be the first index omitted by
$D$, taking $k=n$ if the search fails. Then $D$ contains exactly the
$a_i$ with $i<k$. Indeed an entry above the first omitted one would force
the omitted one into $D$. Conversely every such prefix is downward
closed. For $k<n$ the prefix consists of precisely the predecessors of
$a_k$ and has order type $k$.

For orders $A,B$ of sizes $n,m$ with sorted arrays $a,b$, if $n\le m$,
the map $a_i\mapsto b_i$ is an order-isomorphism onto the prefix of
length $n$. The comparison of $n,m$ therefore decides equality of types
or the direction of the proper-initial-segment relation. Any isomorphism
onto an initial segment maps a point's predecessors bijectively onto all
predecessors of its image. Counting predecessors forces its value on
$a_i$ to be $b_i$. Thus this isomorphism is unique, every automorphism
is the identity, and no finite order is isomorphic to a proper initial
segment of itself. The identity, inverse and composite of finite
isomorphisms are finite maps preserving and reflecting the order. Their
existence makes isomorphism an equivalence relation, and equality of
sizes is its complete decision procedure for finite linear orders.

A strictly increasing sequence of $s$ indices in $I_n$ has its $j$th
entry at least $j$, by induction from one entry to the next. The
descending version gives $d_j+j\le d_0$ for a decreasing sequence of
natural numbers. Thus no descending chain in $A$ starting at $x$ has
more than $r_A(x)$ strict steps. The chain through ranks
$r_A(x),r_A(x)-1,\ldots,0$ attains the bound.

\paragraph{Induction and recursion along a finite presentation.}
For a fixed primitive recursive Boolean predicate $P(x,\vec u)$, suppose
that at each point $x$ the truth of $P$ at all predecessors implies
$P(x,\vec u)$. Induct on the Boolean
$\bigwedge_{i<j}P(a_i,\vec u)$ for $j\le n$. At stage $j$ the
induction hypothesis covers exactly the predecessors of $a_j$.
This proves $P$ throughout the order. If a bounded search has a
counterexample, the least failing rank instead supplies its least
counterexample. Given a fixed primitive recursive step $H$, a finite
recursion with history has an equally explicit realization. Start with
the empty list and append
$H(a_j,(v_0,\ldots,v_{j-1}),\vec u)$ at stage $j$. Primitive
recursion on $j$ builds the whole table and the answer at $x$ is the
entry at $r_A(x)$. Two candidate tables satisfying these equations agree
on their first $j$ entries by induction on that bounded prefix test.
The step $H$ is fixed for each instance of the schema.

\paragraph{Sums, products and powers as order types.}
The concatenation of $A$ and $B$ has consecutive ranks $i$ and $n+j$
on the two blocks. In the ordered family of $m$ copies of $A$ compare
pairs $(j,i)$ by $j$ first, then $i$. The rank is $nj+i$, since $j$
whole blocks precede this copy and $i$ entries precede the chosen entry
within it. For $n>0$, quotient and remainder by $n$ decode this rank. For $n=0$ the product is empty. These descriptions prove that ordinal
sum and product restricted to finite ordinals are natural addition and
multiplication. In particular both are commutative on this restriction.

Give the words of length $m$ over $I_n$ the order comparing the highest
place of disagreement first, and write $d_0$ for the least significant
place. Their ranks are $\sum_{i<m}d_i n^i$. The empty word has rank
zero and is the sole word of length zero, including for $n=0$.
For the next length, fixing the highest digit $j$ produces a copy of
the preceding order, so the block calculation gives rank
$j n^m+r$ and exactly $n^{m+1}$ words. This proves the finite
exponentiation formula and handles the singleton alphabet as well.

\paragraph{Cofinality and topology in the finite case.}
A sublist is cofinal if every point is less than or equal to some entry.
For $n>0$ this says precisely that it contains $n-1$. Thus the least
order type of a cofinal subset is $1$, and for the empty order it is $0$.
In the order topology on a finite chain, every interior singleton is
the interval between its two neighbors. Endpoint singletons are open
rays, and in a one-point order the whole space is open. Expressing a
mask as the union of its occupied singletons proves that every coded
subset is open and closed and that its derived set is empty. The empty
chain also satisfies these statements.

\subsection{Membership as a binary digit}

Interpret $x\in a$ as the assertion that digit $x$ of $a$ in base two
is one. One primitive recursive expression for this digit is
$\lfloor a/2^x\rfloor\bmod 2$. Thus $a$ codes a hereditarily finite set.
Indeed $x\in a$ implies $2^x\le a$, and $x<2^x$ follows by numerical
induction. Consequently $x<a$. Equality of all digits below
$1+\max(a,b)$ implies $a=b$, by the unique binary expansions, which can
also be proved by repeated division with remainder by two. Set inclusion
and the condition $\operatorname{Tr}(a)$ that $a$ be transitive are
therefore bounded predicates. In particular, in $x\in y\in a$ all three
positions are bounded by $a$. Finite unions are computed by bitwise OR.

\paragraph{The finite hierarchy.}
Define $\nu(0)=0$ and
\[
                  \nu(n+1)=\nu(n)+2^{\nu(n)}.
\]
The added bit was zero since $\nu(n)<2^{\nu(n)}$. All lower bits
remain as before, and all higher bits are zero. Thus this is the
successor $a\mapsto a\cup\{a\}$. Its first four values are $0,1,3,11$.
The codes strictly increase, so $i\le\nu(i)$ and the index is unique.
Iterating the one-bit description gives
\[
        x\in\nu(n)\quad\Longleftrightarrow\quad
                     (\exists i<n)\ x=\nu(i).
\]
One can fix $x$ as a parameter in this induction. In particular
$\nu(i)\in\nu(j)$ exactly when $i<j$, and
$\nu(i)\subseteq\nu(j)$ exactly when $i\le j$. For the reverse
inclusion test, if $j<i$, then $\nu(j)$ belongs to $\nu(i)$ and fails
to belong to itself. These formulas prove transitivity of each code,
as well as strict linearity and the least-element property of
membership on its elements. They identify the canonical collapse of an
order $A$ as the map $x\mapsto\nu(r_A(x))$ onto the members of
$\nu(n)$.

\paragraph{Recognition from intrinsic set properties.}
For every number $a$ the following four conditions are equivalent.
\begin{enumerate}
\item $a=\nu(i)$ for some $i\le a$.
\item $a$ is transitive and each of its members is transitive.
\item $a$ is transitive and its members are strictly well-ordered by
membership, with subsets represented by finite masks.
\item $a$ is transitive and membership on its members is a strict
linear order.
\end{enumerate}
We already proved that the first condition gives the next three. The
third gives the fourth. For the fourth implying the second, take
$z\in y\in x\in a$. Set transitivity puts $y,z$ in $a$, so
transitivity of the order there gives $z\in x$. This proves
transitivity of $x$.

For the second implying the first, use numerical course-of-values
induction through $a$. Each member $y$ is transitive, and each member
of $y$ is a member of $a$, hence transitive. The induction gives
$y=\nu(i)$ for some $i\le y<a$. The indices $i$ of the members of $a$
form a downward closed subset of $I_a$, since $j<i$ implies
$\nu(j)\in\nu(i)$ and $a$ is transitive. It is a prefix $I_k$, by the
segment argument above. Thus $a$ and $\nu(k)$ have precisely the same
members and coincide by binary extensionality. Here the induction
invariant can be taken as the conjunction of the claimed bounded
implications for all codes up to the current stage. This is
quantifier-free after replacing bounded scans by their primitive
recursive characteristic functions. An efficient intrinsic test uses
the second or fourth condition. A simpler equivalent decision
procedure checks equality with $\nu(i)$ for $i\le a$.

For a nonempty finite family of such codes, intersection is the member
of least index and union the member of greatest index, by the inclusion
test. These are respectively the infimum and supremum. The empty union
is $\nu(0)$. The successor of the supremum is strictly above every
member of the family. Hence no hereditarily finite set can contain all
the $\nu(i)$ as elements. More generally, for any code $a$ the specific
code $\nu(a+1)>a$ cannot be a member of $a$.

\subsection{The finite Mostowski collapse}

Let $R$ now be any directed relation on $I_n$, and call it extensional
if distinct vertices have different predecessor masks. Acyclicity is
decidable by checking the absence of directed cycles of length $1$
through $n$, since removing the intervening part between repeated
vertices shortens a nonsimple cycle until a simple cycle remains.
Equivalently, every nonempty mask contains a point with no predecessor
in that mask. For, if a mask contained none, repeatedly choose the
numerically least predecessor in the mask. Within $n+1$ visited points
two would coincide, producing a directed cycle. The finite pigeonhole
step follows by counting the distinct occupied bits as above.
Conversely a cycle's vertices give a nonempty mask with a predecessor
at each point.

\paragraph{Collapse theorem.}
An acyclic extensional finite relation is isomorphic to membership on
a unique transitive hereditarily finite set, with a unique isomorphism.
To construct it, repeatedly remove the least numbered point with no
predecessor in the current mask. Its predecessors, if any, have already
been removed. This gives a topological list $v_0,\ldots,v_{n-1}$. At
stage $j$ define
\[
                  c(v_j)=\bigvee_{i<j,\,R(v_i,v_j)} 2^{c(v_i)},
\]
where the join denotes bitwise OR, with empty join zero. After $j$
stages, the earlier codes are distinct. For the induction step suppose
$c(v_j)=c(v_h)$ with $h<j$. Extensionality gives a point in the
symmetric difference of their predecessor masks. It occurs earlier
than $v_j$ in the topological list. Its code contributes a bit on one
side of the displayed equality. Equality supplies an earlier
predecessor with that same code on the other side, and injectivity on
the first $j$ positions identifies the two vertices, a contradiction.
This establishes injectivity throughout. Reading off the bits now gives
$R(x,y)$ exactly when $c(x)\in c(y)$. The range set
$C=\bigvee_{j<n}2^{c(v_j)}$ is transitive because each member of a
$c(v_j)$ is a $c(v_i)$ from an earlier stage. Conversely, let a
bijection $f$ identify $R$ with membership on another transitive
set $D$. Every member of $f(v_j)$ belongs to $D$ by transitivity,
hence is the image of a predecessor. The membership equation for
$f(v_j)$ therefore has precisely the same right-hand side as for
$c(v_j)$, once equality on earlier entries is known. Induction up the
list and binary extensionality give $f=c$, then $D=C$. The proof covers
the empty relation on an empty domain as well.

All loops in the construction have numerical lengths fixed by $n$ and
the already computed table. Values can grow rapidly, but appending the
next value to the current table is a single fixed primitive recursive
operation. This distinguishes primitive recursiveness from a polynomial
bound on output size.

\paragraph{Lean realization.}
The accompanying \texttt{Finite.lean} defines \texttt{Mem}, \texttt{Sub},
\texttt{TransSet}, \texttt{von}, and a bounded search
\texttt{isVon}. Its theorem \texttt{elements} gives the exact bit
description. Theorems \texttt{ordinal\_mem\_iff},
\texttt{ordinal\_sub\_iff}, and \texttt{structural\_recognition}
prove respectively the order tests and the transitive-set recognition test.
The helper \texttt{cut x b} computes one plus the last occupied ordinal
index below $b$, or zero. For a transitive $x$ this gives the prefix
length directly. The same file defines the Boolean order condition \texttt{OrderOn} and computes \texttt{rank} and \texttt{unrank}. Theorems \texttt{rank\_lt\_iff}, \texttt{rank\_unrank}, and \texttt{unrank\_rank} prove the order comparison and both inverse identities. Its proof of surjectivity uses the finite pigeonhole principle, proved there by removing the last input and moving its image into the available output position. To extend the file to the general collapse construction,
encode the relation in digit $ny+x$ of a natural number for the pair
$(x,y)$, and use the digit and bounded-search lemmas for masks.
Implement removal, sorting, topological sorting and collapse by a common
stage variable $j=0,\ldots,n$. The invariants in the three construction
proofs above give their induction motives, with the partially constructed
tables as primitive recursive functions of the fixed input and $j$.

\section{Successors, fundamental sequences and cofinality}
\label{sec:fundseq}

Write $P\mathbin{\|}Q$ for list concatenation, $[b]$ for a one-entry
list, and $[b]^k$ for $k$ copies of $b$ in a list.
All arguments of ordinal comparisons in this section are canonical terms.
We use the convention that a limit ordinal is positive.
The root of even the empty term counts as one node. Thus its node count is
\[
 v([a_1,\ldots,a_k])=1+\sum_{i=1}^k v(a_i).
\]

\begin{proposition}[Successor and finite remainder]
The immediate successor of $x$ is $S(x)=x\mathbin{\|}[0]$.
A nonempty term is a successor exactly when its last exponent is zero.
Every term has a unique expression $S^r(L)$ with $r$ a natural number and
$L$ zero or a positive term ending in a positive exponent.
\end{proposition}
\begin{proof}
A proper extension of $x$ begins with an additional exponent at least zero,
so it is at least $S(x)$. If a term exceeds $x$ by a first unequal
exponent, that same difference puts it above $S(x)$.
Thus $x<y$ is equivalent to $S(x)\le y$.
Deleting a last zero inverts $S$. Counting and deleting all trailing
zeros gives $r,L$ and reversing this operation recovers the input.
Maximality of the trailing zero run proves uniqueness. In particular the
finite terms are the lists $[0]^r$. The terms below $\omega=[1]$, where
$1=[0]$, are exactly these lists, because every positive exponent is at least
$1$. Thus $\omega$ is the first infinite term.
\end{proof}

For a term $\lambda=P\mathbin{\|}[b]$ ending in $b>0$, define
\[
 F(\lambda,n)=
 \begin{cases}
 P\mathbin{\|}[c]^{n+1},& b=S(c),\\
 P\mathbin{\|}[F(b,n)],& b\text{ ends in a positive exponent}.
 \end{cases}
\]
Give $F$ the default value zero on other inputs. The computational call
in the second line follows the last child in the input tree.
The number of calls is bounded by $v(\lambda)$. The two clauses give
$v(F(\lambda,n))\le(n+1)v(\lambda)$ by addition of node counts along
this path. One can therefore implement the definition by bounded
descent with a stack of saved prefixes and a reverse traversal for
reassembly. For instance,
\[
 F(\omega,n)=n+1,\qquad
 F(\omega^2,n)=\omega(n+1),\qquad
 F(\omega^\omega,n)=\omega^{n+1}.
\]
Here numerals in ordinal expressions stand for the corresponding finite
terms. The indexing convention starts each example strictly above zero.

\begin{theorem}[A quantitative cofinality theorem]
\label{thm:fsquant}
If $\lambda$ ends in a positive exponent, $F(\lambda,n)$ is canonical,
and for every $n$,
\[
 F(\lambda,n)<F(\lambda,n+1)<\lambda.
\]
More precisely, whenever $y<\lambda$ and $n\ge v(y)$, one has
$y<F(\lambda,n)$.
\end{theorem}
\begin{proof}
In the successor-exponent clause, every entry of $P$ is at least $b>c$.
The extra copy of $c$ gives strict increase by proper extension. The
first replacement of $b$ by $c$ places both results below $\lambda$.
In the limit-exponent clause the same conclusions hold by applying the
comparisons for the smaller tree $b$ to the only varying entry. This
also justifies canonicality at every return from the recursive call.

For the quantitative claim compare $y$ with $P\mathbin{\|}[b]$. If
$y$ ends in or at $P$ as a prefix, the result has a nonempty suffix
after $P$ and exceeds $y$. If $y$ first differs inside $P$, the
smaller entry of $y$ still decides the comparison. The remaining case is
$y=P\mathbin{\|}[d]\mathbin{\|}Q$ with $d<b$. If $b=S(c)$ then
$d\le c$. For $d<c$ the first new entry decides. For $d=c$ all
entries of $[d]\mathbin{\|}Q$ are at most $c$, and this list has
at most $v(y)$ entries. A list of $n+1$ copies of $c$ therefore
exceeds it at a first smaller entry or by proper extension. If $b$ ends
in a positive exponent, apply the claim to the smaller input pair
$(b,d)$, since $v(d)<v(y)\le n$. It gives $d<F(b,n)$ as required.

These verifications use ordinary induction on the length of the
calculation. Explicitly, for fixed inputs store the visited pairs
$(\lambda,y),(b,d),\ldots$ and the common $n$. At each pair the local
comparison either decides the result outright or invokes the next pair.
Check the implications backwards with a bounded conjunction over the
visited nodes. For strict increase use the analogous list of rightmost
ancestors with parameters $n,n+1$. This makes the induction predicate
a primitive recursive Boolean test.
\end{proof}

\begin{corollary}[The exact first cofinality index]
\label{thm:exact-fs}
Let $M(y)$ be the largest length of a run of identical consecutive
children anywhere in the tree $y$, with empty maximum zero.
For each positive limit $\lambda$ and $y<\lambda$ a primitive recursive
function $m(\lambda,y)$ satisfies
\[
   y<F(\lambda,n)\quad\Longleftrightarrow\quad n\ge m(\lambda,y),
   \qquad m(\lambda,y)\le M(y).
\]
In particular $M(y)$ can replace $v(y)$ in
Theorem~\ref{thm:fsquant}. For every natural $k$, equality holds in
the bound when $\lambda=\omega$ and $y=k$.
\end{corollary}
\begin{proof}
Write $\lambda=P\mathbin{\|}[b]$. If $y$ is a prefix of $P$,
including equality, or first differs from $\lambda$ inside $P$,
return zero. The whole prefix $P$ survives and is followed by at
least one entry, so the desired strict comparison holds at every
index. Otherwise $y=P\mathbin{\|}[d]\mathbin{\|}Q$ with $d<b$.
If $b=S(c)$ and $d<c$, return zero again. If $d=c$, return the length
$r$ of the initial run of $c$ in $[d]\mathbin{\|}Q$. The new suffix
of $F(\lambda,n)$ has $n+1$ copies of $c$. For $n+1\le r$ this suffix
is a prefix of the suffix of $y$, possibly equal. For $n+1>r$ it
exceeds that suffix because $y$ ends after the run or continues with
a smaller entry. This proves the threshold $n\ge r$, and $r\le M(y)$.
If $b$ is a limit, return $m(b,d)$. The new suffix has the single
entry $F(b,n)$, so $y<F(\lambda,n)$ is equivalent to $d<F(b,n)$.
In particular equality at that entry cannot give strict inequality,
even when $Q$ is empty. The recursion descends into proper subtrees, and
$M(d)\le M(y)$ as $d$ is a child of $y$. A bounded stack of these
input pairs and a backwards scan prove both assertions by the same
quantifier-free induction as before. Computing $M$ itself requires
only counting runs at each node and taking their finite maximum.
Finally $F(\omega,n)=n+1$ and $M(k)=k$, so $m(\omega,k)=k$.
\end{proof}

\begin{example}
Let $\lambda=\omega^{\omega^2}$ and $y=\omega^{\omega\cdot3+7}+42$.
Comparison of the leading exponents reduces the first-index problem to
$\omega^2$ and $\omega\cdot3+7$. The latter begins with three copies of
the exponent $1$, so $m(\lambda,y)=3$. Indeed
$F(\lambda,2)=\omega^{\omega\cdot3}<y<F(\lambda,3)=\omega^{\omega\cdot4}$.
The final $42$ units do not enter the calculation.
\end{example}

\begin{corollary}[Characterizations and cofinalities]
For $x>0$ the following four conditions are equivalent.
\begin{enumerate}[nosep,label=(\roman*)]
\item The last exponent of $x$ is positive.
\item The term $x$ has no predecessor.
\item $S(y)<x$ for every $y<x$.
\item The strict predecessors of $x$ have least upper bound $x$.
\end{enumerate}
The cofinalities of zero, a successor and a positive limit in this
presentation are respectively $0,1,\omega$. Here a cofinal list or
sequence below $x$ must supply an entry $\ge y$ for each $y<x$.
\end{corollary}
\begin{proof}
For a positive last exponent, take $n=v(y)$ in
Theorem~\ref{thm:fsquant}. Then $S(y)\le F(x,n)<x$.
Every proposed upper bound $y<x$ of the predecessors is exceeded by
$F(x,n)$. A final zero, on the other hand, exhibits a predecessor $p$
which is itself the greatest strict predecessor and their least upper
bound. This separates the two cases constructively.

The empty list suffices at zero and the singleton containing $p$ suffices
at $S(p)$, where an empty list fails on $p$. At a positive limit the
displayed $F$ gives an increasing sequence. It exceeds the maximum of
each nonempty finite list of predecessors at the explicit index given
above. The empty list misses zero. Since the terms below $\omega$ are
the finite ones, this also proves minimality of the stated cofinality.
In particular $\omega$ is the least positive limit. The three values give
$\operatorname{cf}(\operatorname{cf}(x))=\operatorname{cf}(x)$ by
inspection.
\end{proof}

\begin{proposition}[A minimum problem already at $\omega+1$]
Let $b(n)$ be any specified primitive recursive Boolean predicate,
possibly with numerical parameters. In $[0,\omega]$ form the decidable
inhabited predicate $A_b$ consisting of $\omega$ and exactly those
finite $n$ for which $b(n)=1$. Then $A_b$ has a least member if and
only if
\[
       (\exists n\ b(n)=1)\ \lor\ (\forall n\ b(n)=0).
\]
\end{proposition}
\begin{proof}
Given a least member, decide whether it is finite or equals $\omega$.
A finite least member supplies a positive instance. If the least
member is $\omega$, any $n$ with $b(n)=1$ would be a smaller member.
Boolean decidability now gives $b(n)=0$ at each input. Conversely,
from a witness $n$ scan the $n+1$ finite positions up to it to find
the first positive one. This is least among all finite members and
lies below $\omega$. In the all-zero case $\omega$ itself is least.
Thus the assertion that all such inhabited decidable subsets have a
minimum already entails the limited principle of omniscience for
primitive recursive predicates, understood as a first-order schema.
\end{proof}

\subsection{The tower endpoint and continuity of the principal-power map}
Put $p(x)=[x]$, the notation for $\omega^x$. Singleton comparison gives
$p(x)<p(y)$ exactly when $x<y$. In fact $x<p(x)$ for every finite
term. This holds at zero. If $a$ is the first child of $x$, then
$a<p(a)$ by the same assertion for the smaller tree, and $p(a)\le x$
by the prefix comparison. Thus $a<x$, whence $x<[x]$ by comparison
of the first entries. This argument is a bounded backwards check on the
leftmost branch. At a positive limit $\lambda$, the sequence
$p(F(\lambda,n))=F(p(\lambda),n)$ is cofinal below $p(\lambda)$.
In particular, given $z<p(\lambda)$, the bound $n=v(z)$ supplies an
argument $u=F(\lambda,n)<\lambda$ with $z<p(u)$. This proves
continuity in the increasing, least-upper-bound sense.

Define $T_0=1$, $T_{n+1}=p(T_n)$, and let $h(0)=0$,
$h(x)=1+\max_i h(a_i)$ for a nonempty term with children $a_i$.
The tower is strictly increasing. In fact,
\[
                 x<T_n\quad\Longleftrightarrow\quad h(x)\le n.
\]
To see this, first prove $x<T_{h(x)}$, and, for $h(x)=d>0$,
$T_{d-1}\le x$. At height zero the upper inequality is $0<1$.
At positive height $d$, all children are below $T_{d-1}$ by their
upper inequalities and monotonicity of the tower. Comparing first
entries gives $x<[T_{d-1}]=T_d$. When $d=1$, nonemptiness gives
$x\ge1=T_0$. When $d>1$, a child of height $d-1$ is at least
$T_{d-2}$ by its lower inequality. The first child is at least this
large and therefore $x\ge[T_{d-2}]=T_{d-1}$. Both estimates can be
verified on the finite table of nodes of the given tree, in postorder.
Together they give the asserted equivalence.

Thus the towers are cofinal in the entire term order, with witness
$h(x)$ for a strict upper bound of $x$. There is also an exact
fixed-point formulation using a one-symbol extension. Adjoin a last
point $E$ and set $p(E)=E$. Then $E$ is the least fixed point, since
$x<p(x)$ for every ordinary term. The towers have supremum $E$, and
$p$ is continuous there since, for $x<E$, $p(T_{h(x)})>x$. This models the
first epsilon fixed point at the endpoint of this presentation. The
unextended term order itself has no largest element.

\section{Order topology by explicit neighborhood tests}
\label{sec:ordtop}
All statements about a family of points below refer to its specified
membership predicate. A basic neighborhood of a nonzero point $x$ can
be shrunk to $(b,S(x))$ with $b<x$, since the right-hand endpoint can
be moved to $S(x)$ and zero supplies a possible lower endpoint.
At zero use $[0,1)$. In an initial segment intersect these intervals
with the segment. A point $x>0$ is therefore a limit point of a
predicate $A$ precisely when for every $b<x$ there is a $y$ with
$b<y<x$ and $A(y)$.

Successors $S(p)$ are isolated by $(p,S(S(p)))$, and zero is isolated.
At a positive limit $x$ the sequence $F(x,n)$ gives points approaching
$x$ strictly from below. The intervals $(F(x,n),S(x))$ form a
decreasing countable local basis by Theorem~\ref{thm:fsquant}. As a
result, an initial segment $[0,a)$ is discrete exactly when
$a\le\omega$. For $a\le\omega$ every point is finite. For $a>\omega$
the segment contains $\omega$ and all the finite points approaching it.
On $[0,\omega]$, a fixed decidable membership predicate $U$ describes
an open set exactly when it omits $\omega$ or it contains $\omega$ and
there is an $N$ such that $U(n)$ holds for every finite $n\ge N$.
Indeed, the isolated points impose no condition. A neighborhood of
$\omega$ contains a final tail and every such tail together with
$\omega$ contains a basic neighborhood. Thus in the case containing
$\omega$ the complement can be given by a finite list, obtained by
checking the $N$ positions below the cutoff.

\begin{theorem}[Exact derivative stages from a last exponent]
Write $D_0(x)$ for true, and for $a>0$ set
\[
       D_a(x)\quad\Longleftrightarrow\quad
       x\ne0\ \text{and the last exponent of }x\text{ is at least }a.
\]
For all term indices the following equations hold, where the prime
means taking limit points in the order topology:
\[
 D_{S(a)}=(D_a)',\qquad
 D_\lambda(x)\ \Longleftrightarrow\ (\forall n)\,D_{F(\lambda,n)}(x)
 \ \Longleftrightarrow\
 (\forall a<\lambda)\,D_a(x)\qquad(\lambda\text{ a positive limit}).
\]
Thus a nonzero $x$ with last exponent $e$ belongs at stage $e$ and
disappears at $S(e)$. Zero disappears at stage $1$.
\end{theorem}
\begin{proof}
Consider $x=P\mathbin{\|}[e]$. If $e\le a$, use $(P,S(x))$.
Each point strictly between $P$ and $x$ must extend $P$ with a
first new entry $d<e$. An earlier discrepancy would put it on the
same side of both bounds, and the prefix ending in $e$ already equals
$x$. Its last exponent is at most $d<e\le a$.
When $a=0$ this open interval has no point below $x$. Thus throughout
the case $e\le a$ it contains no other $D_a$ point.
Zero has its own isolating neighborhood.

If $e>a$, fix $b<x$. For $e=S(c)$, the term $F(x,v(b))$ is strictly
between $b$ and $x$ and ends in $c\ge a$. For $e$ a limit, take
$n=\max(v(b),v(a))$. Then $F(x,n)>b$ and its last exponent $F(e,n)$
exceeds $a$, again by Theorem~\ref{thm:fsquant}. This supplies a
$D_a$ point for every prescribed neighborhood and proves the successor
equation.

For the limit equation, $D_\lambda(x)$ implies membership at every
smaller stage by comparison of indices. For nonzero $x$ with last
exponent $e<\lambda$, the one index $F(\lambda,v(e))>e$ witnesses
failure of both universal conditions. For $x=0$, the term
$F(\lambda,0)>0$ witnesses the same failure, since our fundamental
sequences use at least one copy in the bottom clause. These are
explicit witnesses for all the negative cases.
\end{proof}

The family $D$ was defined directly by a primitive recursive test on
two inputs. The displayed equations certify its behavior at every
stage without applying a transfinite recursion principle in iPRA.
Each $D_a$ contains its limit points since $D_{S(a)}\subseteq D_a$.
It is cofinal in the whole term order. For a proposed bound $x$ the
singleton $[S(\max(x,a))]$ exceeds $x$ because $x<[x]$, and its
exponent exceeds $a$. These give explicit closed unbounded families on
this particular countable notation order. The neighborhood proofs also
apply pointwise within $[0,b)$. When $b>0$ the stage $D_b$ in this
segment is empty because, for each positive $x<b$, its last exponent is no
larger than its first, which was proved strictly below $x$.

\section{Controlled descending chains below a fixed finite power of omega}
\label{sec:control}
Fix a positive integer $d$. Terms below $\omega^d=[d]$ are exactly
the polynomials in $\omega$ with natural coefficients and exponents
$d-1,\ldots,0$ (zero coefficients fill the gaps). Indeed, the first
exponent, if present, must be $<d$, and every smaller entry is then
also a finite exponent $<d$. Reading multiplicities from highest to
lowest identifies the term order with strict lexicographic order of
$d$-tuples of naturals. At the first unequal coefficient, the tuple
with the shorter run either ends or passes to a lower exponent, which
proves the comparison rule in both directions.

Let $g(i)$ be a given nondecreasing primitive recursive bound on
all coordinates allowed at time $i$. Define an end index for each
fixed dimension by
\[
 H_0(s)=s+1,\qquad
 H_{j+1}(s)= H_j^{\,g(s)+1}(s),
\]
where the superscript on $H_j$ means ordinary iteration of that
natural-number function. Each $H_j$ is nondecreasing and
$H_j(s)\ge s+1$. Indeed, increasing the
start or the number of iterations cannot decrease the result, and
there is at least one iteration. For every fixed $j$ the iterate
uses only primitive recursion.

\begin{theorem}[A length bound with a computable failure index]
A list occupying times $s,\ldots,t-1$ of strictly decreasing
$j$-tuples satisfying the coordinate bound has $t\le H_j(s)$
(for $s\le t$). In particular every primitive recursive sequence of
terms below $\omega^d$, for fixed $d$, has a primitive recursively
located index $i$ with $f(i+1)\ge f(i)$.
\end{theorem}
\begin{proof}
At dimension zero there is just the empty tuple, so the list has
at most one entry. Suppose the bound is proved in dimension $j$.
In a nonempty decreasing list of $(j+1)$-tuples the first coordinates
are nonincreasing and start at most $g(s)$. Split into maximal
constant-first-coordinate blocks by scanning the list. There are at
most $g(s)+1$ blocks because every new block drops the first coordinate
by at least one. If a block starts at $u$, its remaining $j$ coordinates
decrease strictly, so its exclusive end is at most $H_j(u)$.
Induction on the number of blocks and monotonicity show that the final
end is at most $H_j^{g(s)+1}(s)$. An empty list satisfies the bound
immediately. The dimension induction constructs a separate iPRA proof
for each fixed $j$. Within each proof the block endpoints are bounded
scans of the supplied list, and the induction on block count only
compares natural numbers computed from those endpoints.

For the last assertion choose $g(i)$ as the maximum of the coordinates
of $f(0),\ldots,f(i)$ (and zero). Apply the bound to the list of
$H_d(0)+1$ consecutive inputs starting at zero. If all adjacent
comparisons decrease, its end index violates the bound. A bounded
scan returns a failing comparison before index $H_d(0)$. Its correctness follows by the same finite argument. The statement for
sequences is a scheme, with $f$ given by a primitive recursive
definition. The dimension is fixed before forming the bound.
\end{proof}

\section{Ordinary arithmetic on hereditary Cantor forms}\label{sec:ordinary}
Write $W(e)=[e]$, $0=[]$, $1=W(0)$, $\omega=W(1)$.
A natural number $n$ in an ordinal expression denotes the list of $n$ zeros.
The length $|x|$ is its number of outer entries, and $h(x)$ is its first
entry when $x>0$. All variables for ordinals in this section range over
valid finite notations. Thus, for example, a formula about all $e$ is an
arithmetical assertion about codes with a decidable validity guard.
Juxtaposition of ordinal variables denotes multiplication once it has been
defined. We first use concatenation of \emph{lists}, written $U\mathbin{\|}V$.

\subsection{Addition, subtraction on the left, and absorption}
Given $a,b$, put $a+0=a$. For $b>0$ retain from $a$ its longest prefix
with entries $\ge h(b)$ and concatenate $b$. This is the normal form of
ordinal addition. To see what happens with many summands, allow an arbitrary
finite word of exponents and scan it from left to right. Each input letter
$s$ removes the current final letters $<s$, then joins the end of the word.
Write $N$ for the result. In any surrounding word, replacing adjacent letters
$r,s$ with $s$ when $r<s$ preserves $N$. Indeed, inserting $r$ can only
have removed letters smaller than $r$. The subsequent insertion of $s$
would remove all of them, as well as $r$. The normalized stack after $s$
is identical in the two scans. The scan itself implements just such
deletions. It follows that $N(U\mathbin{\|}V)$ is unchanged by normalizing
either factor first. For descending words it is exactly the displayed
prefix rule.

\begin{theorem}\label{th:ordinary-add}
Addition is associative and has two-sided unit $0$. For fixed $a$, the
map $b\mapsto a+b$ is an increasing bijection onto the terms $\ge a$,
with a primitive recursive inverse on its range. For fixed $b$ the map
$a\mapsto a+b$ is weakly increasing. In particular, left cancellation
holds, $a+(b+1)=(a+b)+1$, and $a+1$ appends one zero.
\end{theorem}
\begin{proof}
Associativity follows by normalizing the concatenation of three words,
and the empty word gives the unit. When a positive summand $d$ removes
a suffix of $x$, its first entry exceeds the first removed entry.
When it removes nothing, $x$ is a proper prefix of $x+d$. Thus $x<x+d$.

For $b<c$ discard their longest common prefix and take $d$ to be the
remaining suffix of $c$. If $b$ ends at the common prefix, $b+d=c$ by
concatenation. Otherwise its first remaining entry is smaller than the
first of $d$. Every remaining entry of $b$ is removed and the common
prefix is retained. Again $b+d=c$. For $b=c$ take $d=0$. Hence
$a+c=(a+b)+d>a+b$ whenever $b<c$. This gives uniqueness of the residual
and the stated range.

For weak monotonicity in the first argument it suffices to insert one
letter $s$ at a time. Deleting the entries below $s$ from two ordered
words preserves their lexicographic order because a first mismatch either
survives, becomes a prefix comparison, or disappears in both.
Appending $s$ to both truncated words preserves that order. Indeed, in
the proper-prefix case the next letter of the longer word is at least
$s$. Equality of this next letter with $s$ leaves the shorter result a
proper prefix. An induction on the scanned length proves the assertion.
\end{proof}

\begin{theorem}[Additive indecomposability]\label{th:add-indecomp}
For $x>0$, the singletons $W(e)$ are exactly the terms closed below themselves under sums of two smaller terms. Equivalently, $x$ is a singleton if and only if it cannot be expressed as a sum of two smaller terms. The map $e\mapsto W(e)$ gives
an increasing enumeration of exactly these terms. Also $u+W(e)=W(e)$
whenever $u<W(e)$.
\end{theorem}
\begin{proof}
A nonempty word is below $[e]$ exactly when its first entry is below
$e$. The sum of two such words still uses only entries below $e$, and
insertion of $e$ deletes all of the first word. The empty-word cases follow
from the unit. If $x$ has at least two entries, split off its final
singleton. Both pieces are smaller than $x$ and concatenation is their
ordinal sum. Singleton comparison compares their sole entries.
\end{proof}

\subsection{Multiplication by expanding the right operand}
Set $0b=a0=0$. For $a>0$ replace each entry $e$ of $b$ by the block
\[
 B_a(0)=a,\qquad B_a(e)=W\bigl(h(a)+e\bigr)\quad(e>0),
\]
and add the blocks in their original order. In particular $aW(e)=B_a(e)$.
For $r<s$, all exponents of $B_a(r)$ are smaller than the sole exponent
of $B_a(s)$. For $r=0$ this uses $h(a)<h(a)+s$. For $r>0$ it uses strict
increase of $h(a)+{}\cdot$. Consequently a deletion $r,s\rightsquigarrow s$
in a word of exponents preserves the sum of its expanded blocks.
The first block has first exponent at least as large as every later
block exponent, so it survives all later insertions. Thus
\[
                  h(ab)=h(a)+h(b)\qquad(a,b>0).                 \tag{1}\label{eq:arith-head}
\]

\begin{theorem}\label{th:ordinary-mul}
Multiplication is associative, has unit $1$ and absorbing element $0$,
and satisfies $a(b+c)=ab+ac$. For $a>0$, $b\mapsto ab$ is strictly
increasing and left cancellation holds. Multiplication is weakly increasing
in its first argument. In particular $a(b+1)=ab+a$.
\end{theorem}
\begin{proof}
Expansion of a concatenation gives the sum of its two expansions.
The preceding deletion argument allows the concatenation to be normalized
first, proving the distributive law. Both unit identities follow from
$B_a(0)=a$ and $B_1(e)=W(e)$ for every $e$. Associativity reduces by
finite distributivity to $a(bW(e))=(ab)W(e)$ for each entry $e$.
Zero factors are immediate and $e=0$ is the unit identity.
For $a,b,e>0$ the two sides are monomials with exponents
$h(a)+(h(b)+e)$ and $(h(a)+h(b))+e$, by (1). They agree by associativity
of addition. If $b<c$, take $d>0$ with $b+d=c$. For $a>0$, expansion
gives $ad>0$, so $ac=ab+ad>ab$. If $0<a\le a'$, then $h(a)\le h(a')$. Each expanded block for $a$ is bounded by its counterpart for $a'$.
Addition is weakly increasing in both arguments, so the same holds for
the sum of all blocks. The case $a=0$ follows from leastness of zero.
\end{proof}

For $s,t\ge0$ this calculation gives
$W(s)W(t)=W(s+t)$, including $t=0$. For a positive $a$ it gives
$a\omega=W(h(a)+1)$. If the first exponent $g$ of $a$ has multiplicity
$j$, then $an$, for a positive finite $n$, has $j(n-1)$ copies of $g$
followed by $a$. Each new copy of $a$ deletes exactly the part of the
preceding copy below $g$. This proves the formula by ordinary induction.
In particular, finite ordinals have their usual addition and multiplication.

\begin{corollary}[Multiplicative indecomposability]
The nonzero terms $x$ satisfying
$u,v<x\Rightarrow uv<x$ are exactly $1$, $2$ and the $W(W(e))$.
\end{corollary}
\begin{proof}
The finite cases $1,2$ can be checked directly. For finite $n\ge3$,
$(n-1)^2\ge n$. Below $W(W(e))$ the heads of nonzero factors are
$<W(e)$, and their sum remains $<W(e)$ by Theorem~\ref{th:add-indecomp}.
Equation (1) bounds their product. Conversely an infinite $x$ with at
least two entries and head $g$ strictly exceeds $W(g)$. Both $W(g)$
and the natural number $|x|+1$ are below $x$, but their product, a list
of $|x|+1$ copies of $g$, exceeds $x$. For infinite $x=W(g)$, if
$g$ is not a singleton then $g>0$ splits as $p+q$ with $p,q<g$.
The two smaller monomials $W(p),W(q)$ multiply to $x$.
\end{proof}

\subsection{Exponentiation with arbitrary base}
The following decomposition is also an indexing theorem for the limits.
Every term $b$ has a unique expression
\[
                         b=\omega Q+n,\qquad n<\omega.            \tag{2}\label{eq:arith-omega-div}
\]
Count its final zeros to obtain $n$. For each remaining entry $e>0$,
solve $1+d=e$ by subtracting one if $e$ is finite, and by taking $d=e$
if $e$ is infinite. Indeed, in the latter case the head of $e$ is
positive, so $1$ is absorbed. The solution is unique and ordered as
$e$ varies by Theorem~\ref{th:ordinary-add}. These solutions in their
given order form $Q$. The rules for multiplication replace its entries
$d$ by $1+d$, also when $d=0$. The recovered positive entries of $b$
and its $n$ zeros prove (2) and uniqueness. Thus $Q\mapsto\omega Q$ is
an order isomorphism onto zero together with the nonzero limits (the
nonempty terms with no final zero).

Write $P_a(n)$ for the ordinary finite product of $n$ copies of $a$,
with $P_a(0)=1$. Set $a^0=1$, $0^b=0$ for $b>0$, and $1^b=1$.
For $a\ge2$ define the code $k_a$ and the general power by
\[
 k_a=\begin{cases}1,&a<\omega,\\ h(a)\omega,&a\ge\omega,\end{cases}
 \qquad
 a^{\omega Q+n}=W(k_a Q)\,P_a(n).                              \tag{3}\label{eq:arith-power}
\]
These operations are primitive recursive. All word scans, prefix cuts and
extractions in (1)--(3) are bounded by the input word length. The recursion
$P_a(n+1)=P_a(n)a$ uses the already defined multiplication. A changing
accumulator is allowed in ordinary primitive recursion. All further scans
are on outputs of earlier primitive recursive operations. For instance,
when $a\ge\omega$, $h(a)>0$ and $k_a=W(h(h(a))+1)$ is a single monomial.
Formula (3) therefore also covers bases and exponents of arbitrarily
large finite nesting depth.

\begin{theorem}[Power laws]\label{th:ordinary-pow}
For all terms $a,b,c$,
\[
 a^{b+c}=a^b a^c,\qquad (a^b)^c=a^{bc},\qquad
 a^{b+1}=a^b a,\qquad \omega^b=W(b).
\]
For $a\ge2$ the exponential is strictly increasing in $b$. For fixed
$b$ it is weakly increasing in $a$.
\end{theorem}
\begin{proof}
First fix $a\ge2$ and put $k=k_a$. By (1) and finite induction,
$h(P_a(m))=h(a)m$, including $m=0$. This head is $<k$. For a finite
base it is zero. For an infinite base and $m>0$ it has head $h(h(a))$, below
the exponent of $k$. In consequence, for $C>0$,
\[
             P_a(m)W(kC)=W(kC).                                \tag{4}
\]
Indeed $kC\ge k$, since multiplication by positive $k$ is increasing.
As $k$ is a single monomial, every exponent of the word $h(P_a(m))$
is smaller than its exponent. More explicitly, if $h(P_a(m))>0$, then
$h(h(P_a(m)))<h(k)\le h(kC)$.
Thus $h(P_a(m))+kC=kC$, proving (4) by the block rule.

Write $b=\omega B+m$, $c=\omega C+n$. If $C=0$, appending $c$ changes
the finite tail of $b$ to $m+n$ and $P_a(m+n)=P_a(m)P_a(n)$ by finite
induction. If $C>0$, $m+\omega C=\omega C$ by absorption, so the sum
has quotient $B+C$ and tail $n$. Removing the middle factors by (4)
in $a^b a^c$ leaves
$W(kB)W(kC)P_a(n)=W(k(B+C))P_a(n)$. This is the first power law.
It includes the successor identity because $a^1=a$.

For the second power law take $b>0$ and define
\[
 D=1\quad(b<\omega),\qquad D=B\omega\quad(b=\omega B+m,\ B>0).
\]
We claim $b\omega=\omega D$ and $k_{a^b}=kD$. Here $a^b\ge2$ because positive
finite powers of $a$ are at least $a$, and $B>0$ gives first exponent
$\ge kB>0$ by (1). If $b$ is finite then $b\omega=\omega$.
A positive finite power of a finite base is finite. For an infinite base,
$h(P_a(b))=h(a)b$ has the same head as $h(a)$, so
$h(P_a(b))\omega=h(a)\omega$. This proves both claims in the finite case.
If $B>0$, equation (1) gives $h(b)=1+h(B)$, since adding the finite
tail preserves the head. Now $b\omega=W((1+h(B))+1)=\omega(B\omega)$.
The head $H=h(a^b)=kB+h(P_a(m))$ has the same head as $kB$ because the second
summand is $<k$, whereas $kB\ge k$, with $k$ a monomial.
Thus $H\omega=(kB)\omega=kD$, proving the second claim.

Apply the first power law repeatedly to obtain
$P_{a^b}(n)=a^{bn}$ by finite induction on $n$.
Since $bc=b\omega C+bn=\omega(DC)+bn$, the first law and (3) give
\[
       a^{bc}=W(kDC)P_{a^b}(n)=(a^b)^c .
\]
If $b=0$, both sides of the second law are $1$. For bases $0$ and $1$
the laws follow by splitting zero from positive exponents and using that
a sum is zero exactly when both summands are zero, and a product is zero
exactly when a factor is zero.

For $a=\omega$, $k=\omega$ and $P_\omega(n)=W(n)$ by the monomial
product formula. Thus (3) is $W(\omega Q)W(n)=W(b)$. For $b<c$, write
$c=b+d$ with $d>0$. Since $a^d>1$ and $a^b>0$ for $a\ge2$, strict
increase of multiplication in the right argument gives $a^c>a^b$.
Finally $k_a$ is weakly increasing for $a\ge2$, by monotonicity of the
head and of multiplication, and $P_a(n)$ is weakly increasing in $a$
by finite induction. Formula (3) proves weak monotonicity in the base.
The exceptional bases satisfy it directly.
\end{proof}

\subsection{Continuity with explicit witnesses}
For a strictly increasing function $F$ on the codes and a nonzero limit
$L$, the limit equation $F(L)=\sup_{d<L}F(d)$ has the following
witnessed meaning. Every $F(d)$ with $d<L$ lies below $F(L)$, and for
every $z<F(L)$ we compute $d<L$ with $z<F(d)$. These two assertions
imply the least-upper-bound property by the decidable linear order.

\begin{theorem}\label{th:ordinary-cont}
The maps $d\mapsto a+d$, $d\mapsto ad$ for $a>0$, and $d\mapsto a^d$
for $a\ge2$ satisfy these limit equations at every coded nonzero limit.
\end{theorem}
\begin{proof}
Strict increase proves the first assertion. For addition, given
$z<a+L$, use $d=0$ if $z<a$. Otherwise compute $q$ with $a+q=z$. Then $q<L$, and $d=q+1<L$ works since $L$ is a limit.
For multiplication, use the Euclidean division theorem for the codes
proved in this paper. Write $z=aq+r$, $r<a$. Then $aq\le z<aL$ gives
$q<L$, and $z<a(q+1)$ with $q+1<L$.

For powers, (2) writes the limit as $L=\omega Q$ with $Q>0$. Thus
$a^L=W(kQ)$ where $k=k_a$. If $z=0$, use $d=0$. If $z>0$, let
$t,r$ be the Euclidean quotient and remainder of $h(z)$ on division
on the left by $k$. Then $h(z)=kt+r<kQ$, $r<k$, and $t<Q$.
If $a$ is infinite, write $g=h(h(a))$, so $k=W(g+1)$. Put
$n=|r|+1$. Every exponent of $r$ is $\le g$, and $h(a)n$ begins with
at least $n$ copies of $g$. Thus $h(a)n>r$. For $d=\omega t+n$
equation (1) now gives
\[
 h(a^d)=kt+h(P_a(n))=kt+h(a)n>kt+r=h(z).
\]
If $2\le a<\omega$, then $k=1$, $r=0$, and $t=h(z)$. Use
$n=|z|+1$. The word $a^{\omega t+n}=W(t)P_a(n)$ consists of
at least $n$ copies of $t$, since the natural number $a^n\ge n$.
It exceeds $z$ by length or at the first smaller exponent of $z$.
In either case $\omega t+n<\omega(t+1)\le\omega Q$. This follows
from $n<\omega$ and left distributivity. All three witnesses have
been constructed by primitive recursive operations.
\end{proof}

For $a=1$ the power function is constant $1$, and for base $0$ the
positive powers are constant $0$. Including the exponent $0$ makes the
supremum for this base equal to $1$, since $0^0=1$. Taking the supremum over
$0<d<L$ gives the article's limit equation in these cases too. The zero
multiplication map has constant value $0$ and also preserves these suprema.

\paragraph{Calculations.}
We have $1+\omega=\omega<\omega+1$ and
$2\omega=\omega<\omega 2$. Thus right cancellation for addition and
multiplication can fail. Also $(1+1)\omega=\omega$ while
$1\omega+1\omega=\omega 2$, exhibiting failure of the other distributive
law. Formula (3) yields $2^\omega=\omega$, $2^{\omega+1}=\omega 2$,
and $\omega^2=W(2)>\omega$. In particular, exponentiation also
fails to commute. With the left argument varying up to $\omega$, the
values $n+1$, $n2$ and $n^2$ on finite arguments are cofinal in
$\omega$, whereas $\omega+1$, $\omega 2$ and $\omega^2$ all exceed it.
This proves failure of continuity in the other argument in all three cases.

The spiral in the opening illustration of \emph{Ordinal number} uses
$f(x)=\omega(1+x)$. Its least fixed point is $W(\omega)=\omega^\omega$.
Indeed, $f(x)>x$ for finite $x$. For $\omega\le x<W(\omega)$ one has
$1+x=x$ and the positive finite head $h(x)$ increases to $1+h(x)$
under multiplication by $\omega$ on the left. At $x=W(\omega)$ the
head becomes $1+\omega=\omega$ and the product is exactly $W(\omega)$.

\subsection{Prime factors and a unique ordered factorization}
An ordinal prime is a term $x>1$ with no expression $x=ab$ with $a,b<x$.
This condition concerns equality. Closure of an initial segment under
all products, treated above, is stronger.

\begin{theorem}[Prime classification and normal factorization]\label{th:ord-primes}
The primes of this notation system are the finite primes, the limit primes
$W(W(d))$, and the infinite successor primes $W(d)+1$ with $d>0$.
Every positive term has a unique prime factorization subject to the
following ordering rules. All limit factors precede all other factors,
the limit factors are nonincreasing, and each uninterrupted group of
finite primes is nonincreasing. All the factors can be computed
primitive recursively.
\end{theorem}
\begin{proof}
If $ab$ is positive and finite, both factors are finite by (1) or
by monotonicity in each argument. The finite primes are therefore
exactly those of elementary arithmetic. The closure assertion above
already proves primality of $W(W(d))$. Let $x=W(d)+1$ with $d>0$
and suppose $x=ab$ with both factors smaller. Each factor must exceed
$1$. If $b$ is a limit, every block of $ab$ is a singleton with
positive exponent, whereas $x$ ends in zero. If $b$ is finite,
$b\ge2$. Then $a$ must be infinite and the formula for $ab$ has at
least two copies of its leading exponent, whereas $x$ has exactly one.
If $b$ is an infinite successor, $ab$ has at least one positive
singleton block before the last copy of $a$. The singleton survives
because its exponent strictly exceeds $h(a)$. As $a>1$ and $ab$ ends in
zero, $a$ ends in zero too and has at least two entries. Thus $ab$ has
at least three entries, whereas $x$ has two.

We now give a factorization algorithm. Combine adjacent equal entries
of $x$ into runs with exponents $e_0<\cdots<e_s$ listed from the
\emph{right} and multiplicities $m_0,\ldots,m_s>0$. For each $i>0$
compute $d_i>0$ with $e_{i-1}+d_i=e_i$ by the prefix algorithm. Write
the word of $e_0$ as $[g_1,\ldots,g_j]$. The ordered factorization,
temporarily keeping finite factors unexpanded, is
\[
 \left(\prod_{t=1}^{j}W(W(g_t))\right)m_0
       \prod_{i=1}^{s}\bigl((W(d_i)+1)m_i\bigr),               \tag{5}
\]
with products read from left to right. The first product is $W(e_0)$
by the monomial product formula and equals $1$ if $e_0=0$.
Multiplication by $m_0$ produces the bottom run. If the portion already
formed is $v$ with head $e_{i-1}$, then
$v(W(d_i)+1)=W(e_i)+v$. Its leading exponent occurs once, so multiplying
by $m_i$ yields exactly $m_i$ copies of $e_i$ followed by $v$.
Finite induction proves (5). Replace each $m_i$ by its finite prime
factors in nonincreasing order, omitting $1$.

Conversely, any product obeying the ordering rules has an initial
block of limit primes, followed by finite blocks separated by the
successor factors $W(d_i)+1$. Finite blocks are allowed to be empty.
Multiplying exactly as above recovers the run exponents and
multiplicities. The bottom exponent records the initial block
letter for letter. Each next exponent determines $d_i$ uniquely by
left cancellation. Each multiplicity recovers its finite block by
unique natural factorization. This proves uniqueness even when a
multiplicity is one.

To see that no further primes occur, take a product (5) containing at
least two prime factors. Its proper prefix before the last factor is
smaller than $x$ by strict increase of multiplication in the right
argument. Its last prime is also smaller. For a finite last prime this
is immediate if $x$ is infinite, and is ordinary arithmetic otherwise.
If the last factor is a limit, the entire list consists of at least two
nonincreasing limit primes, so the bottom exponent $e_0$ exceeds its
last singleton $W(g_j)$. If the last factor is $W(d)+1$, the preceding
product $v$ exceeds $1$ and $x=W(h(v)+d)+v$. Here $h(v)+d\ge d$. Even in the case of equal leading exponents, the tail $v$ exceeds $1$.
Thus (5) splits into two strictly smaller factors. A single prime
factor and the empty product account for all remaining cases.

For the elementary arithmetic used here, factor a natural $m>1$ by
trial division up to $m$, removing a least divisor greater than one
at each stage. Each such divisor is prime by minimality, and at most
$m$ nonterminal divisions occur since the remaining quotient decreases.
Uniqueness uses Euclid's lemma. If a prime $p$ does not divide $v$, the
ordinary Euclidean algorithm on $(p,v)$ returns gcd $1$, since the only
positive divisors of $p$ are $1,p$. Back-substitution gives integers
$u,w$ with $up+wv=1$. Initialize the two inputs by their two unit
coefficient vectors and replace a row by the difference of that row
and the quotient times the other. Multiplication by $z$ shows that
$p\mid vz$ implies $p\mid z$. Euclid's lemma locates an equal prime
in any second factor list, after which cancellation and finite induction
prove uniqueness. The decreasing Euclidean remainders permit the fixed
fuel $p+v+2$. All scans, products, and sorting steps in (5) are bounded
iterations on finite codes.
\end{proof}

\section{Division and the order types represented by the operations}
\label{sec:division}
Write $A+B$, $AB$ for the ordinary ordinal operations just constructed. Within a CNF, $\|$ means list concatenation. A sum may delete final entries from its left input.
For $X\le Y$ denote by $Y\mathbin{\dot-}_{\!L}X$ the suffix of $Y$ after its longest common prefix with $X$. Theorem~\ref{th:ordinary-add} proves that this is the unique $D$ with $X+D=Y$. We will apply this inverse both to exponents and to full terms.
\begin{theorem}[Euclidean division of ordinal notations]\label{thm:orddiv}
For $A>0$ and every $B$, a PR algorithm computes the unique pair $(Q,R)$ with
\[
 B=AQ+R,\qquad R<A .
\]
\end{theorem}
\begin{proof}
Let $d$ be the first exponent of $A$. Split $B=H\|V$ after all exponents $e>d$. All exponents in $V$ are at most $d$. For each entry $e$ in $H$ solve $d+u=e$ by left subtraction. Every $u$ is positive. Since $d+(-)$ is strictly increasing, these $u$ form a weakly descending list $U$. The multiplication rule gives $AU=H$ entry by entry, also when $U$ is empty.

Let $m$ be the number of initial entries $d$ of $V$ (zero if there are none). If $A$ has $k\ge1$ leading entries $d$, the finite multiple $A j$, for $j>0$, has exactly $kj$ such leading entries followed by the tail of $A$ below $d$. In particular $A(m+1)>V$, even if $m=0$. Among $0,\ldots,m$ select the greatest $j$ with $Aj\le V$. Zero qualifies, and
$Aj\le V<A(j+1)$. Solve $Aj+R=V$ by the residual algorithm. If $R\ge A$, monotonicity in the right argument gives $V\ge Aj+A=A(j+1)$, a contradiction. Return $Q=U\|[0]^j$. Left distributivity, associativity and the canonical boundary between $H$ and $V$ give
\[
 AQ+R=(AU+Aj)+R=H+V=B.
\]
If $Q<Q'$ then $AQ+R<A(Q+1)\le AQ'\le AQ'+R'$ for $R<A$. The reversed inequality is symmetric. Thus equal results force equal quotients and, by left cancellation of addition, equal remainders. This also proves the maximality characterization $AQ\le B<A(Q+1)$. Each search in the construction has a stated finite bound.
\end{proof}
\begin{example}
Take $A=\omega+3$ and $B=\omega^2+\omega\cdot2+1$. The high block gives $U=\omega$. In the tail, $A\cdot2=\omega\cdot2+3$ is already too large. The algorithm chooses $j=1$ and returns $Q=\omega+1$, $R=\omega+1$. Adding $R$ erases the final three units of $A$. Ordinary natural-number division of the two leading coefficients alone would miss this case.
\end{example}
\begin{corollary}[Ordered sum, product, and intervals]
The tagged sum of $I_A$ followed by $I_B$ is order-isomorphic to $I_{A+B}$. The lexicographically ordered pairs $(j,i)$ with $j<B$, $i<A$, with the block index $j$ compared first, are order-isomorphic to $I_{AB}$. Both isomorphisms and their inverses are PR on the specified domains. In particular $I_B$ is isomorphic to the interval $[A,A+B)$.
\end{corollary}
\begin{proof}
Send the first summand identically to $I_A$ and the second by $y\mapsto A+y$. To invert, decide $x<A$ and otherwise compute its left residual from $A$. Strict increase shows that the residual is $<B$. For products send $(j,i)$ to $Aj+i$. Each block lies between $Aj$ and $A(j+1)$, with the latter endpoint excluded. The map is strictly increasing within and across blocks. For $x<AB$ its quotient by $A$ must be $<B$, so division inverts the map. When $A=0$ or $B=0$ both sides of the product claim are empty. These arguments require neither a least element of an arbitrary subset nor a general order-type assignment.
\end{proof}

\begin{theorem}[Cardinality of each represented initial segment]
If $a\ge\omega$, there is a uniform PR bijection, with PR inverse, between $\mathbb N$ and the canonical numerical codes in $I_a$. The least infinite notation of that cardinality is $\omega$. Distinct order types of this cardinality include $\omega$ and $\omega+1$.
\end{theorem}
\begin{proof}
Let $p_a(c)$ test that $c$ is a valid notation strictly below $a$. The $n+1$ finite terms $0,\ldots,n$ all pass and have distinct codes. Compute $B(n)$ as the maximum of their numerical codes, including the code of zero. Scan the integers up to $B(n)$ and return the $(n+1)$st passing integer $e_a(n)$, in \emph{numerical code order}. Counting the marked positions among the $n+1$ distinct witnesses guarantees that the scan has at least that many hits. For an accepted code $c$, set $r_a(c)=\sum_{i<c}p_a(i)$. The unit increments of this counter imply $r_a(e_a(n))=n$. Conversely if $r_a(c)=n$, then $c$ is the unique passing integer of rank $n$, hence $e_a(n)=c$. Bounded-sum invariants prove both identities. The same enumeration works with validity alone, so all the canonical names are effectively countable. An injection from $\mathbb N$ into a finite segment would restrict to one more point than its size, contradicting finite counting. Below $\omega$ lie just the finite terms. The segment at $\omega+1$ has a greatest point $\omega$, whereas each point $n$ below $\omega$ has the larger point $n+1$ there, excluding an order isomorphism. At the same time $\omega\cdot\omega=\omega^2>\omega$ and the resulting two segments are in PR bijection.
\end{proof}

\subsection{Exponentiation as an order on finitely supported functions}
\label{sec:exporder}
Fix canonical terms $a\ge2$ and $b$. A function from $I_b$ to $I_a$
of finite support is represented by a list of pairs $(s,d)$ in strictly
decreasing order of $s<b$, with $0<d<a$. At every unlisted argument its
value is zero. Order two distinct lists by their digits at the largest
argument where they differ. That argument is found by merging their
two finite lists of occupied positions.

\begin{theorem}[The order-type interpretation of ordinal powers]
Uniformly in $a,b$, the map
\[
                f\longmapsto V(f)=\sum_{s\text{ in decreasing order}}
                                      a^s f(s)
\]
is a primitive recursive order isomorphism onto $I_{a^b}$, with a
primitive recursive inverse. This includes the empty support and $b=0$.
\end{theorem}
\begin{proof}
First allow arbitrary finite lists of occupied positions, without an
upper bound $b$. For a nonempty list with highest position $s$ its
value lies in $[a^s,a^{s+1})$. To prove this, consider its first
digit $d$ and the value $L$ of the tail. If the tail is empty,
$L=0<a^s$. Otherwise its highest position $t<s$ gives
$L<a^{t+1}\le a^s$ by induction on the finite list length. Thus
\[
 a^s\le a^s d\le a^s d+L<a^s d+a^s=a^s(d+1)\le a^{s+1}.
\]
Here $d+1\le a$ follows from the successor property. Likewise the
tail strictly below any position $s$ has value less than $a^s$.
At the greatest differing position write the two digits as $d<d'$,
allowing zero. The smaller digit and its tail then contribute strictly
less than $a^s(d+1)\le a^s d'$. Restoring the identical higher terms
by addition on the left preserves the strict inequality. This proves
order preservation and reflection, and gives injectivity. The same
bounds show that $V(f)<a^b$ exactly when all occupied positions are
below $b$, including the empty list since $a^b>0$.

It remains to produce the digits of any given $x$. Put $k=k_a$ as in
\eqref{eq:arith-power}, so $a^\omega=W(k)$ and
$a^{\omega J+i}=W(kJ)a^i$ for finite $i$. For each outer exponent $e$
of $x$, divide on the left by $k$, writing $e=kJ+r$ with $r<k$.
The intervals $[kJ,k(J+1))$ are ordered by $J$, and inside such an
interval the order is the order of its remainders. Consecutive entries
with a common $J$ therefore give a canonical word $C_J$ of their
remainders, including repeated entries and zeros. Its head is $<k$, so
$0<C_J<W(k)$. Each original group has value $W(kJ)C_J$, since
$W(kJ)W(r)=W(kJ+r)$, including $r=0$. Summing these groups in
decreasing order of $J$ reconstructs $x$.

Find a natural $N$ with $C_J<a^N$ as follows. For finite $a$, $k=1$,
so all entries of $C_J$ are zero. Its value is its length $m$,
and $N=m+1$ gives $a^N>m$. For infinite $a$ put $g=h(h(a))$.
Then $k=W(g+1)$. Every entry of the head $r_*=h(C_J)$ is at most
$g$. With $N=|r_*|+1$, ordinary finite multiplication shows that
$h(a)N$ begins with at least $N$ copies of $g$. Thus
$h(a)N>r_*$, by the first differing exponent or by length when all
of $r_*$ consists of copies of $g$. Since $h(a^N)=h(a)N$,
again $C_J<a^N$.

Now set $q_0=C_J$. For $i<N$ divide $q_i=aq_{i+1}+d_i$,
where $d_i<a$. From $q_i<a^{N-i}$ one gets
$q_{i+1}<a^{N-i-1}$, since $aq_{i+1}\le q_i$ and
$a\,a^{N-i-1}=a^{N-i}$. At the last step $q_N<1$, hence $q_N=0$.
Substituting the equations back, using left distributivity at each
step, gives $C_J=a^{N-1}d_{N-1}+\cdots+a^0d_0$.
Multiplication on the left by $W(kJ)$ moves these digits to the positions
$\omega J+i$. These positions increase with $i$, and positions from
distinct groups are ordered by $J$, since
$\omega J+i<\omega(J+1)$. Omitting zero digits and concatenating the
groups therefore returns a valid finite support. It has value $x$.
The initial-segment bound already proved guarantees the required
positions $<b$ whenever $x<a^b$. Uniqueness follows from the strict
comparison proved above.

There are at most $|x|$ groups, and each group has its displayed
numerical bound $N$ on the division loop. Intermediate lists are
computed by ordinary primitive recursion with a changing accumulator.
This construction uses no induction on the ordinal quotient.
\end{proof}
For $a=1$ there is only the zero-valued function, and $1^b=1$.
For $a=0$ there is one empty function when $b=0$ and none when $b>0$,
in agreement with the exceptional powers. In this last case a total
function on a nonempty domain cannot choose a value at its least input.

\section{Natural arithmetic, indecomposability, and finite nimbers}
\label{sec:natural}
Write $c_a(e)$ for the number of occurrences of $e$ among the children of $a$. To compare coefficients, scan the finite union of supports in decreasing order. At the largest exponent with different coefficients the term with smaller coefficient is smaller. This follows by expanding equal leading runs and comparing the next entry, including the case of a missing entry.

\begin{theorem}[Hessenberg operations]\label{thm:hess}
Define $a\mathbin\oplus b$ by merging the two lists in decreasing order. Define $a\mathbin\otimes b$ by listing all $e\oplus f$ with $e$ an entry of $a$, $f$ an entry of $b$, with their repetitions, and sorting in decreasing order. These PR operations make $\mathcal O$ a commutative semiring. Natural sum strictly preserves order in each argument, as does natural product when the other factor is positive. For finite inputs they give ordinary natural sum and product.
\end{theorem}
\begin{proof}
Multiplicity in a merge is $c_a(e)+c_b(e)$. Induction on the entries inserted proves this for each queried exponent. Equality of all coefficients on the union of supports implies equality of the sorted lists by matching equal runs. Thus $\oplus$ is commutative and associative with identity $0$. It preserves the largest coefficient difference and its sign. In the product the coefficient at $g$ is the finite sum of $c_a(e)c_b(f)$ over support pairs with $e\oplus f=g$. Swapping the pair gives commutativity. Triples $(e,f,h)$ contribute equal numbers to the same exponent under either bracketing because $(e\oplus f)\oplus h=e\oplus(f\oplus h)$. Splitting a merge in the second factor splits each coefficient into two contributions and proves distributivity. The empty list has no pairs and $1=[0]$ contributes one pair with each left entry, preserving it. Lists of zeros give the finite laws.

For strict product monotonicity take $a<b$ and $z>0$. Let $e$ be the largest exponent at which $a,b$ differ, and $f$ the leading exponent of $z$. Above $e\oplus f$ every contribution from the first factor comes from an exponent greater than $e$, and the coefficients there agree. At $e\oplus f$ those larger exponents still contribute equally. Smaller first exponents cannot contribute. With first exponent $e$ only second exponent $f$ contributes. Strict increase of $\oplus$ proves each exclusion. The excess at this exponent is $(c_b(e)-c_a(e))c_z(f)>0$ as required. All supports and the pair and triple indexing sets are finite, bounded by products of the input lengths.
\end{proof}

\begin{proposition}[Comparison with ordinary arithmetic]
$a+b\le a\oplus b$ and $ab\le a\otimes b$. Repeated $n$-fold natural sum of $a$ equals $a\otimes n$. The two natural operations can fail to preserve suprema in either argument.
\end{proposition}
\begin{proof}
The sum algorithm deletes entries, whereas the merge keeps every entry. Deletion from a decreasing list can only decrease its value. For the product first take a single right monomial $[f]$. For $f=0$ both results are $a$. For $f>0$ and $a>0$ the ordinary result is $[d+f]$, with $d$ the leading exponent of $a$. The natural result has leading exponent $d\oplus f\ge d+f$, with at least one copy. For an arbitrary right list fold these ordinary products in order, using the bound for a sum and monotonicity of natural sum at each step. Distributivity reconstructs the natural product. Zero causes no exception. The identity for $n$ follows from the pair formula, or ordinary induction. Finally $1\oplus n=n+1$ has supremum $\omega$, but $1\oplus\omega=\omega+1$. Likewise $2\otimes n=2n$ has supremum $\omega$, but $2\otimes\omega=\omega\cdot 2$. For either sequence every value is finite and an index exceeding any proposed finite bound can be given by its successor.
\end{proof}

\begin{corollary}[Closure at principal terms]
The terms below $\omega^t$ are closed under natural sum. Those below $\omega^{\omega^t}$ are closed under natural product.
\end{corollary}
\begin{proof}
Every entry of either summand below $[t]$ is $<t$, hence this remains so after a merge. For positive factors below $[\,[t]\,]$, let their leading exponents be $e,f<[t]$. The largest product exponent is $e\oplus f<[t]$ by the first assertion. Zero factors give zero directly.
\end{proof}

\paragraph{A polynomial reading.}
Natural arithmetic here is the semiring of finite polynomials with natural coefficients in commuting variables $X_t$ indexed by valid terms $t$. Assign $X_t=[\,[t]\,]$. A monomial $\prod_t X_t^{m_t}$, using natural product and its finite powers, has one exponent, the decreasing list with $m_t$ copies of $t$. Natural sum assembles such monomials as the children of a term. Conversely, sort the children and their own children to read off uniquely the monomials and their multiplicities. The two conversions are inverse and the pair formula gives multiplication of polynomials. This is an equality of finite coded data. The universal polynomial evaluation property for any fixed coded commutative semiring follows by evaluating this finite list with its given operations.

\subsection{The finite nimber sum}
The article also links ordinals as nimbers. For natural numbers let $a\mathbin{\mathrm{xor}}b$ be bitwise exclusive-or. For a finite list of integers write $\mathrm{mex}$ for its least missing natural. A list of length $k$ misses some value among $0,\ldots,k$ by counting, giving a bounded algorithm.
\begin{theorem}
For finite nimbers the mex addition rule has the unique solution $a\boxplus b=a\mathbin{\mathrm{xor}}b$.
\[
 a\boxplus b=\mathrm{mex}\big(\{i\boxplus b:i<a\}\cup\{a\boxplus j:j<b\}\big).
\]
In particular this operation is commutative and associative, has identity zero, and $a\boxplus a=0$.
\end{theorem}
\begin{proof}
Put $z=a\mathbin{\mathrm{xor}}b$. Exclusive-or by a fixed input is its own inverse, so no legal option gives $z$. For $w<z$, find the highest bit at which $w$ and $z$ differ by a bounded scan. At that bit $z$ has one and $w$ zero. Exactly one of $a,b$ has one there. If it is $a$, take $i=w\mathbin{\mathrm{xor}}b$. All higher bits of $i$ and $a$ agree, and this bit decreases, giving $i<a$ and $i\mathbin{\mathrm{xor}}b=w$. The other case yields $j=w\mathbin{\mathrm{xor}}a<b$. Thus the first missing value is precisely $z$. The other assertions follow bit by bit. For uniqueness, tabulate the rule in any fixed rectangle $i\le a,j\le b$ by increasing $i+j$. Agreement with exclusive-or for all earlier entries is a bounded invariant and forces agreement at the next entries. In particular $1\boxplus1=0$, while $1+1=1\oplus1=2$.
\end{proof}

\section{An extension through the first epsilon numbers}
\label{sec:epsilon}
A finite name for an exponentiation fixed point need not unfold into a tree of infinite depth. We now change the alphabet, retaining the old terms as an initial segment. This gives a constructive counterpart of the article's enumeration $\varepsilon_0,\varepsilon_1,\ldots$ of fixed points. All the structural operations in this section, including the enumeration and its cofinality witnesses, are uniform in the numerical index.

\paragraph{Syntax and comparison.}
A word is a finite weakly decreasing list of atoms. An atom is either a marker $M_i$ for $i\in\mathbb N$, or $P(u)$ with $u$ a word. Put $E_i=[M_i]$. Forbid $P(E_i)$ in a valid word. Thus $E_i$ already serves as its own power of $\omega$. Give a marker $M_i$ level $i+1$. The level of $P(u)$ is the level of $u$. The level of a word is the largest level among its atoms, with empty maximum zero. Compare atoms by level first. Within a positive level the marker is least, and two $P$ atoms compare by their arguments. In level zero all atoms are $P$ atoms. Compare words lexicographically as before. The left and right inputs being finite tagged trees, scanning and recursion on their numerical subcodes give PR tests. The mutual transitivity proof is the same first-difference proof as for Theorem~\ref{thm:order}. At the atom stage use transitivity of levels, then the least-marker case, then the smaller pair of argument trees. For a precise induction order give each word and each atom a tagged-node code with every child smaller. Use the bounded triple invariant of that theorem for both node sorts. This proves trichotomy, equality reflection and transitivity on valid codes.

For nonempty words the leading atom has maximal level. Hence every word of smaller level precedes every word of larger level, and $E_i$ is the least word of level $i+1$. On words of level zero erasing each $P$ tag gives exactly the original CNF order, with inverse obtained by retagging every exponent. In particular those words form $I_{E_0}$. Successor appends $P(0)$, the least atom, by the same prefix argument as before. Define
\[
 W(u)=\begin{cases}E_i,&u=E_i,\\ [P(u)],&\text{otherwise}.\end{cases}
\]
It preserves level. Across levels it therefore strictly preserves order. Within one level its ordinary case uses argument comparison. When exactly one argument is a marker singleton it is least in that level both before and after applying $W$. Thus $W$ is strictly increasing.

\begin{theorem}[Exactly the named fixed points]
For every extended word $u$, $u\le W(u)$, with equality exactly for the $E_i$. The map $i\mapsto E_i$ enumerates them in strictly increasing order, and every word is below some $E_i$.
\end{theorem}
\begin{proof}
The empty word is below $[P(0)]$. At a marker singleton the result is the definition. For a different nonempty $u$ let $A$ be its first atom. If $A$ is a marker, $A<P(u)$ by their common level and the marker rule. If $A=P(v)$, validity excludes a marker singleton $v$. Structural induction gives $v<W(v)=[A]\le u$. Comparing the two $P$ atoms gives $P(v)<P(u)$. In both cases $u<[P(u)]$. This induction uses a finite table of claims on subcodes. The markers have distinct levels in the given order. A word of level $k$ is strictly below $E_k$, which has level $k+1$.
\end{proof}

\begin{theorem}[Cofinal witnesses below each marker]
Set $B_0=1$ and $B_{i+1}=S(E_i)$. For each $i$, the strict sequence $T_i(n)=W^n(B_i)$ has supremum $E_i$. On input $y<E_i$, a number of iterations bounded by the syntax depth of $y$ already gives a strict upper bound for $y$. The sequence of the markers itself has no upper bound among the extended words.
\end{theorem}
\begin{proof}
$B_i$ is a non-marker singleton or a proper extension of a marker singleton. Its level is $i$, and all iterates are non-marker words of the same level. Hence they strictly increase below $E_i$. Define $d(0)=0$, and for a nonempty word set $d(y)=1+\max\big(\{0\}\cup\{d(v):P(v)\text{ occurs as an outer atom of }y\}\big)$. Prove $y<T_i(d(y))$ for $y<E_i$. Empty is strictly below $B_i$. For a nonempty $y$ of depth $h$ compare its head with the head of $T_i(h)=[P(T_i(h-1))]$. A marker head below $E_i$ has level at most $i$ and is strictly smaller than this $P$ atom. A head $P(v)$ has $v<E_i$ and $d(v)\le h-1$. Structural induction and increase of the iterates give $v<T_i(h-1)$. If the atom levels differ, the smaller level settles comparison, and at equal level the last strict inequality settles it. This proves the claim by induction over the finite tree, allowing $i$ as a numerical parameter. For the least-upper-bound assertion every smaller candidate is passed by the computed iterate. Finally $E_{\mathrm{level}(y)}$ exceeds any given word $y$.
\end{proof}
Externally, interpret $M_i$ as $\varepsilon_i$ and $P(u)$ as $\omega^{|u|}$. Every word containing markers no larger than $M_j$ denotes an ordinal below $\varepsilon_{j+1}$, by structural induction and closure of that fixed point under finite CNF sums and powers with smaller exponents. A word whose largest marker is $M_j$ is at least $\varepsilon_j$ by induction and $\xi\le\omega^\xi$. At a common level $P(u)$ with a marker somewhere and $u\ne E_j$ lies strictly above its level's marker. To see this, prove simultaneously by tree height that the only term of level $j+1$ denoting $\varepsilon_j$ is $E_j$. A $P(v)$ at that level has $v\ne E_j$ by validity and the induction hypothesis gives $|v|>\varepsilon_j$. Standard CNF comparison now identifies the displayed comparison with the classical one. Conversely below $\varepsilon_0$ the original terms suffice. Between $\varepsilon_j$ and $\varepsilon_{j+1}$ the towers from $\varepsilon_j+1$ are cofinal by the usual continuity of $\omega^{(-)}$ and the definition of the next fixed point. All ordinals below their base have names by induction on $j$, including the marker itself. Passing from one tower bound to the next, the exponents in the CNF of any ordinal below the new bound are below the previous bound. They already have names. Use the marker atom for a named fixed exponent, and $P$ for the others. This proves that the extended words name precisely the ordinals below $\varepsilon_\omega$. In the arithmetic, adding a fresh greatest symbol $E_\omega$ makes it the supremum of the markers, and extending $W(E_\omega)=E_\omega$ preserves strict increase. This final extension is only asserted for the order and $W$.

\section{Cuts, normal functions, and the reach of induction}
\label{sec:scope}
In this section $\mathcal O$ again means the original system below $\varepsilon_0$.
\subsection{Von Neumann identities as formulas about points}
Represent a possibly infinite collection by a given membership formula. Write $I_a(x)$ for $x<a$ with the usual guards. Then
\[
 I_a\subseteq I_b\ \longleftrightarrow\ a\le b,\quad
 I_a\cap I_b=I_{\min(a,b)},\quad I_{S(a)}=I_a\cup\{a\}.
\]
Here inclusion means pointwise implication. Transitivity proves inclusion when $a\le b$. If $b<a$, the point $b$ refutes it. The other formulas follow from comparison and immediacy of the successor. Thus even infinite represented ordinals admit a precise membership and inclusion reading in the arithmetic. For $a>0$, the pointwise union of its predecessors' cuts is $I_a$ if $a$ is a limit, and $I_c$ if $a=S(c)$. In the first case, $x<a$ implies $x<S(x)<a$, which gives an explicit predecessor containing $x$. In the second case $I_c$ itself occurs in the union and contains all its other summands. The empty case has empty union.

For a fixed PR sequence of codes $b(n)$, its union cut has the first-order membership formula $\exists n\ (x<b(n))$. If all $b(n)\le L$ and each $x<L$ has a witness $n$ with $x<b(n)$, the upper bound and the witness give the two inclusions proving that the union equals $I_L$. In general a formula describing a union need not yield a computable canonical name for it. For example let $b_e(n)$ be $1$ if a particular machine $e$ has halted by step $n$, and $0$ otherwise. These increasing-in-the-weak-sense sequences are uniformly PR. A total algorithm computing a canonical supremum for all of them would decide whether $e$ halts by testing the result against zero. If it halts the supremum is $1$, and if it never halts the supremum is zero. This is an external computability obstruction, already for a two-valued sequence.

\subsection{Two consequences of cofinality}
A fixed PR function $f$ on notations is called normal here if it is strictly increasing and has the following limit property. Whenever $z<f(L)$ at a positive limit, some specified or existentially given $y<L$ satisfies $z<f(y)$. For fixed $f$ its iterates $t_0=s$, $t_{n+1}=f(t_n)$ are PR. Suppose a positive limit $L$ is supplied with $t_n<L$ for all $n$ and with $\forall x<L\ \exists n\ x<t_n$. Then $L$ is the least fixed point of $f$ at or above $s$. Indeed $f(t_n)<f(L)$, so the tail sequence lies below $f(L)$. To see that the tail is cofinal, use $t_0<L$ and cofinality to get a term strictly above $t_0$, whose index is positive, and replace any witness at $t_0$ by that term. Thus $L\le f(L)$. If $L<f(L)$, continuity gives $y<L$ with $L<f(y)$. Choose $n$ with $y<t_n$. Then $f(y)<t_{n+1}<L$, a contradiction. Decidable comparison gives equality. For any fixed point $p\ge s$, numerical induction on the comparison $t_n\le p$ keeps all iterates below or equal to $p$, and so $L\le p$. The fixed-point argument uses the given limit name and cofinality. The towers supply both in the one-symbol extension with endpoint $E$.

Closed unbounded subsets of the whole original term order need not have nonempty intersection. In the tower sequence $T_0=1$, $T_{n+1}=[T_n]$, take separately the even and the odd indices. Both ranges are cofinal by the height bound and strict growth. Both membership tests are bounded. For $x$ of tree height $h$ only indices $n\le h+1$ could have $T_n\le x$. Each range has no limit point \emph{inside} $\mathcal O$. For $x>0$ choose the greatest listed point of that parity below $x$, when present. The interval $(p,S(x))$ then excludes all listed points except possibly $x$. If none is present, the open ray below $S(x)$ serves. Zero is isolated. The ranges are disjoint by strictness. This illustrates the uncountable regularity hypothesis in the article's club-intersection theorem.

There is also a useful cover statement at the first limit. Suppose a cover of $[0,\omega]$ has supplied choices of a covering open set at every point, by a PR selector (including an index for $\omega$), and the neighborhood data for the set chosen at $\omega$ give a tail bound $N$. Keep that one set and the choices at $0,\ldots,N-1$. Comparison of a point with $N$ proves that the list covers, using the tail for the remaining finite points and for $\omega$. This proves the finite-subcover form with the indicated data. Together with the neighborhood description it gives the witnessed one-point compactification of the discrete naturals.

\subsection{Hereditary bases and a sharp descent boundary}
For $b\ge2$ call a valid term $b$-admissible if every equal-exponent run, at every node, has length $<b$. Evaluate a word at base $b$ by
\[
 V_b([e_1,\ldots,e_k])=\sum_{j=1}^k b^{V_b(e_j)}.
\]
Define $H_b(n)$ by taking the base-$b$ digits of $n$ and using $H_b(j)$ as the exponent at each nonzero digit position $j$, with that digit as multiplicity. The expansion of zero is empty. When $n>0$, every used position $j$ has $j<n$ because $b^j\ge j+1$. Thus $H_b$ is defined by a course-of-values table on $n$, uniformly in $b$. Digits are found by bounded division and remainder, and the evaluation is a table on subtrees with $b$ a parameter.

\begin{theorem}[Base change and strict local descent]
For each $b\ge2$, $H_b$ and $V_b$ are inverse, order-preserving bijections between the naturals and the $b$-admissible terms. If $c\ge b$ and $n>0$, put $m=V_c(H_b(n))$. Then $m>0$ and
\[
       H_c(m)=H_b(n),\qquad H_c(m-1)<H_b(n).
\]
\end{theorem}
\begin{proof}
For admissible words, evaluation reflects the first CNF coefficient difference. To see this without assuming admissibility of any output, first prove this comparison lemma for pairs of admissible input trees. At each node the corresponding comparison assertion for child exponents is available by smaller-subtree induction. Apply it also to pairs of children from the same input, so distinct exponent values become distinct integer positions in the same order. If $j$ is the greatest position with a differing coefficient, the excess at $j$ is at least $b^j$ and the entire smaller-position part on either side is at most $\sum_{i<j}(b-1)b^i=b^j-1$. The last identity is ordinary induction. The empty-difference case gives equal values. This proves comparison including equality. For the induction one may fix the finite list of all subterms of the inputs and prove the comparison test for all pairs in it with maximum height at most $s$.

Now prove simultaneously by numerical table induction up to $n$ that $H_b(i)$ is valid and admissible and $V_b(H_b(i))=i$ for all $i\le n$. Distinct earlier digit positions have correct integer values, so their names are correctly ordered by the comparison lemma. Digits supply the permitted multiplicities, and ordinary base expansion proves the evaluation identity. For an admissible $u$, the trees $H_b(V_b(u))$ and $u$ have equal evaluations and hence are identical by equality reflection. These calculations give both inverse and order assertions. Finally $b$-admissibility implies $c$-admissibility and a nonempty word evaluates positively. Apply the inverse identity at $c$ and strict comparison to $m-1<m$.
\end{proof}
A special Goodstein step consists precisely of changing $b$ to $b+1$ and subtracting one from the positive value. The theorem proves decrease for each step and, by finite iteration, checks every supplied finite trace up to its first zero. Universal termination of all such sequences is unprovable even in PA, which interprets iPRA with the same PR definitions. This is the classical independence result of Kirby and Paris. See Rathjen's proof in the sources. Accordingly the stepwise decrease does not furnish a uniform termination proof in iPRA. The bound for every fixed $\omega^d$ in Section~\ref{sec:control} is a positive contrast with an explicit failure index.

\subsection{Why arbitrary presentations are different}
Here is an obstruction adapted from ProveIt's two-type promise note. Insert the number $0$ into the usual order of the positive integers by placing $x>0$ below it precisely when machine $e$ has not halted by time $x$. Positive-positive comparisons are unchanged. The positives placed below zero form an initial segment by monotonicity of bounded simulation, so we have a uniformly PR strict linear order. If $e$ never halts its order type is $\omega+1$. If $e$ first halts at $t$, zero comes immediately before $\max(1,t)$ in the positive list. The order type is $\omega$. A computable normalizer for presentations even with this promise, returning names with decidable equality to $\omega$, would decide the halting problem. This external argument explains why the algorithms take canonical notation, or explicit finite presentations with certificates. The article's proper class of all ordinals, unrestricted transfinite induction or recursion, uncountable cardinals, and stationary sets on regular uncountables have no literal quantifiers in the first-order numerical language. Finite-list escape ($S(\max A)$ above every listed term), witnessed suprema, the fixed-dimensional descent theorem and the two club families above are precise arithmetical counterparts.

\section{Lean formalization and proof obligations}
\label{sec:formal}
Lean 4.32.0 checked the following theorems with its standard library. The object-language construction in the next subsection gives the remaining steps for certification in iPRA.

\subsection{What the kernel has checked}
\renewcommand{\arraystretch}{1.15}
\begin{tabularx}{\linewidth}{@{}p{.22\linewidth}X@{}}
\toprule
File & Definitions and proved results\\\midrule
\texttt{Finite.lean} & Ackermann bit membership, exact elements of $\nu(n)$, membership and inclusion tests, and intrinsic recognition by two levels of transitivity. On any Boolean finite linear order, the predecessor count and its bounded inverse are order-preserving bijections. See \texttt{structural\_recognition}, \texttt{rank\_lt\_iff}, \texttt{rank\_unrank} and \texttt{unrank\_rank}.\\
\texttt{OrdinalCore.lean} & Binary-tree syntax, recursive comparison, validity, $+$ and $\cdot$. Equality reflection, reversal, transitivity, comparison of principal terms, strict successor growth and agreement of addition on finite inputs.\\
\texttt{Limits.lean} & Immediate successor, height of canonical terms and the exact tower test, classification by the last exponent, and fundamental sequences. \texttt{fundamental\_sequence} proves validity, increase and the limit bound. \texttt{cofinal\_modulus} proves the node-count witness.\\
\texttt{Threshold.lean} & \texttt{firstIndex\_spec} proves the exact threshold equivalence for a bounded search. The search uses the node-count bound. The sharper maximum-coefficient bound and the direct recursive formula are proved in Corollary~\ref{thm:exact-fs}.\\
\texttt{Power.lean} & General addition associativity, two inverse equations for division by $\omega$ with validity of the quotient, and general power as an executable definition. Finite-exponent recursion and transfinite examples are checked.\\
\bottomrule
\end{tabularx}

Here is the mathematical content of the two main limit declarations, with Boolean tests written as predicates. For valid $x$ ending in a positive exponent, $F(x,n)$ is valid and $F(x,n)<F(x,n+1)<x$. For valid $y<x$ and $v(y)\le n$, the comparison $y<F(x,n)$ holds. These use the same convention $F(\omega,n)=n+1$ as the manuscript. The Lean tree is either \texttt{nil} or \texttt{cons head tail}, so a word and a suffix are both trees. In the node count a cons contributes the count of its head plus that of its tail, and nil contributes one. This is exactly $1+\sum v(e_i)$.

\subsection{From the definitions to object-language derivations}
The ProveIt snapshot used here defines \texttt{LeanProofs.ConciseIPRA.PR k}, \texttt{PTm k}, \texttt{Fm PTm k} and \texttt{Prf PRAAx}. A declaration with conclusion
\texttt{([] : List (Fm PTm k))} $\vdash$\texttt{[PRAAx]} $\varphi$ represents an object-language derivation with $k$ free numerical variables. Universal introduction can close those variables. Its numeric primitives include composition and primitive recursion. The finite algorithms in this report therefore specify their symbols by explicit PR definitions, without adding axioms about their intended computations.

\begin{enumerate}[leftmargin=*,label=\arabic*.]
\item \textbf{Arithmetic representation.}
Use \texttt{IPRADerivations.Core.Definitions} for \texttt{PTm.toPR}, \texttt{recDef}, \texttt{rec\_zero}, and \texttt{rec\_succ}. Use \texttt{Core.Pairing} for the square pairing, inverse projections, and paired registers, and \texttt{Core.Search} or \texttt{Core.SearchParameters} for least bounded search. Encode cons by one plus the pairing from Section~\ref{sec:foundation}. Set $\mathrm{entry}(L,0)=\mathrm{head}(L)$ and $\mathrm{entry}(L,i+1)=\mathrm{entry}(\mathrm{tail}(L),i)$, with default zero. Implement retrieval by iterating tail then taking head. Prepend each new result to a history list, retrieving the entry for stage $j<t$ at index $t-1-j$. The defining equations and paired-register projections give the readback lemma by numerical induction with the queried $j$ fixed. A subcode table is an instance of this history. On malformed values all tests and constructors have fixed default outputs, and theorem hypotheses guard semantic conclusions.
\item \textbf{Bounded induction and order.}
Introduce Boolean operators, a bounded conjunction and a first-failure witness as in Section~\ref{sec:foundation}. To prove an instance $A(s)$ of a bounded universal assertion, apply the repository's \texttt{ind} rule to the equality $\mathrm{allPass}(s)=1$. Eliminate bounded conjunctions by instantiation at a selected index. Introduce them by a scan, retaining a failure witness if the result is zero. Build $\mathrm{cmp}$ from the empty, head and tail clauses. To avoid a choice of traversal order in a rectangular table for comparison, increase the first input's numerical code $i$ and at each stage fill every second input $j\le B$, for fixed bound $B$. Both recursive comparisons have smaller first code and bounded second code. Derive comparison transitivity by the bounded triple invariant of Theorem~\ref{thm:order}. Reflect equality of codes using injectivity of the pairing. The checked Lean proof \texttt{cmp\_trans} supplies the exact branch split.
\item \textbf{Finite structures.}
Represent a relation by its $n^2$ membership bits, a partial table by a numerical list, and an order by its bounded validity test. Sorted-prefix invariants certify rank/unrank and isomorphism. Vertex-removal invariants certify the collapse. For $\nu$, prove first the one-bit successor identity using \texttt{Core.Powers} and division with remainder. The bounded recognition formula checks both levels of transitivity by ranging over indices $<a$. Use the finite table of all implications up to the current numerical code for strong induction. The semantic Lean \texttt{Sub} contains a universal quantifier, but the arithmetic induction test uses only the finite digit ranges.
\item \textbf{Limit sequences, derivatives and descent.}
Define the tagged classifier to return zero, a predecessor, or a value $F(x,n)$ as in \texttt{Limits.localStep}. The paper's rightmost-child recursion also gives a PR definition by tables. For the modulus introduce PR symbols $V,L,C,F$ for validity, ends-in-positive, strict comparison and the fundamental sequence, and $v$ for node count. The object theorem has free inputs $x,y,n$ with matrix
\[
 (V(x)=1\ \wedge\ V(y)=1\ \wedge\ L(x)=1\ \wedge\ C(y,x)=1\ \wedge\ v(y)\le n)
 \ \longrightarrow\ C(y,F(x,n))=1 .
\]
Here $u\le n$ is the quantifier-free equality $u\mathbin{\dot-}n=0$ on natural numbers. Fix $x,y,n$ and tabulate the implications for all subcodes of $x,y$ by decreasing recursive dependency. Every recursive pair is within the original boxes, and node counts of relevant children decrease. The \emph{parameter} $n$ is fixed. For consecutive-index growth, compute the two values at $n$ and $n+1$ in parallel. For derivatives define $D_a(x)$ by its last-exponent test, prove the local positive/negative neighborhood witnesses, and only then introduce the universal quantifiers of the topology statements. For fixed $d$, construct $H_d$ at the Lean meta-level by $d$ repetitions of the iterator construction. The proof by blocks is a bounded induction on the input list.
\item \textbf{Arithmetic and order models.}
Represent a word-normalization scan by its time and stack. Prove equality of stacks immediately after consuming an absorbable pair. Use this lemma before the algebra of addition, then the block-expansion lemma before multiplication. General division uses only these two operations and the bounded maximum in Section~\ref{sec:division}. Powers need the split at $\omega$, finite powers, and the leading-exponent identity. After addition and multiplication, construct division by $\omega$ and then general division. Next prove the power laws and continuity, which supply the bounds needed for digit expansions. For rank/unrank of function words, check digit bounds and carry the finite high-to-low prefix sum as accumulator. For prime factorization use runs and ordinary bounded integer division. The inverse procedure in equation (5) proves uniqueness. The object-level proofs of the general division, exponent, and prime-factor theorems remain to be implemented.
\item \textbf{Natural arithmetic and extensions.}
Sort and count by bounded insertion. Prove multiplicity invariance on the finite union of input and output supports. Products use a rectangular table of pairs, and associativity reindexes triples. The module \texttt{Core.Fubini} provides a reusable finite-sum exchange. For nimbers implement digits and exclusive-or by a bounded sum of digit differences modulo two. For mex use first missing below one plus list length. For the epsilon extension use a tag for marker versus power, encode its index or exponent as a component, and compute levels before comparisons. The three inputs to a transitivity proof shrink together on recursive comparisons. Track the maximum input level and a syntax-height cap for the tower theorem. For hereditary base change use a table through the input integer and a second for evaluating the resulting tree. Digit comparison needs a conjunction over all pairs of subterms including pairs from the same root. These finite invariants suffice without any induction on the notation order.
\end{enumerate}
The normality and cover results are conditional theorems or schemata over fixed algorithms. Introduce their witnesses before discharging hypotheses and closing the formulas. A completed object derivation can then be transported by the repository's \texttt{pra\_iff\_native} equivalence. This object-language stage remains open for the collection.

\section{Source correspondence and further questions}
\small
\begin{description}[leftmargin=.8em,labelindent=0pt,style=nextline]
\item[Order and sets, Sections~\ref{sec:foundation}, \ref{sec:finite-models}, \ref{sec:division}, \ref{sec:scope}.] Strict order, order types, initial segments, finite Mostowski collapse, von Neumann membership, union, and cardinality. \href{https://en.wikipedia.org/wiki/Ordinal_number}{Ordinal number},
\href{https://en.wikipedia.org/wiki/Well-order}{Well-order},
\href{https://en.wikipedia.org/wiki/Order_type}{Order type},
\href{https://en.wikipedia.org/wiki/Mostowski_collapse_lemma}{Mostowski collapse lemma},
\href{https://en.wikipedia.org/wiki/Cardinal_number}{Cardinal number}.
\item[Limits and topology, Sections~\ref{sec:fundseq}--\ref{sec:control} and \ref{sec:scope}.] Successor, cofinality, order topology, induction, least points, and the club-intersection hypothesis. \href{https://en.wikipedia.org/wiki/Successor_ordinal}{Successor ordinal},
\href{https://en.wikipedia.org/wiki/Limit_ordinal}{Limit ordinal},
\href{https://en.wikipedia.org/wiki/Cofinality}{Cofinality},
\href{https://en.wikipedia.org/wiki/Order_topology}{Order topology},
\href{https://en.wikipedia.org/wiki/Transfinite_induction}{Transfinite induction},
\href{https://en.wikipedia.org/wiki/Club_set}{Club set}.
\item[Arithmetic and fixed points, Sections~\ref{sec:ordinary}--\ref{sec:epsilon}.] Ordinary and Hessenberg operations, order-type models, prime factorization, finite mex arithmetic, normal functions and epsilon numbers.
\href{https://en.wikipedia.org/wiki/Ordinal_arithmetic}{Ordinal arithmetic}, including its Cantor normal form and prime factorization sections,
\href{https://en.wikipedia.org/wiki/Additively_indecomposable_ordinal}{Additively indecomposable ordinal},
\href{https://en.wikipedia.org/wiki/Nimber}{Nimber},
\href{https://en.wikipedia.org/wiki/Normal_function}{Normal function},
\href{https://en.wikipedia.org/wiki/Epsilon_numbers_(mathematics)}{Epsilon numbers}.
\item[Constructive and proof-theoretic context.] N. Kraus, F. Nordvall Forsberg and C. Xu, \emph{Type-Theoretic Approaches to Ordinals}, TCS 957 (2023), 113843, surveys constructive tradeoffs and develops binary-tree CNFs in type theory. M. Rathjen, \emph{Goodstein's theorem revisited}, gives the metatheoretic nonprovability used in Section~\ref{sec:scope}. Our local descent proof uses finite digit arithmetic. \href{https://arxiv.org/abs/2208.03844}{Kraus--Nordvall Forsberg--Xu},
\href{https://arxiv.org/abs/1405.4484}{Rathjen}.
\item[ProveIt.] We used the theory definition and derivation interface from Vladimir Reshetnikov's repository, at commit \texttt{95aa4ea543d15ec628998c1f21cd71147745728e}. Section~\ref{sec:scope}'s two-type example is adapted from the repository note on effective ordinal semantics. The proof packet is standalone.
\href{https://github.com/VladimirReshetnikov/ProveIt/tree/95aa4ea543d15ec628998c1f21cd71147745728e/Logic/PrimitiveRecursiveArithmetic}{iPRA directory},
\href{https://github.com/VladimirReshetnikov/ProveIt/tree/95aa4ea543d15ec628998c1f21cd71147745728e/Logic/PrimitiveRecursiveArithmetic/Derivations/Lean/IPRADerivations/Core}{Core derivations},
\href{https://github.com/VladimirReshetnikov/ProveIt/blob/95aa4ea543d15ec628998c1f21cd71147745728e/Computability/HilbertTenthProblem/Papers/research-wip/native-stream-queue/ordinal_two_type_effectivity_boundary.md}{Two-type promise note}.
\end{description}
\paragraph{Further questions.}
Can the numerical code for a CNF be replaced by a word code with smaller readback bounds while shortening the object derivation of quantitative cofinality? What uniform bounds can be proved when both the descent dimension and the growth-control function are restricted by explicit input data? For the extended epsilon notation, can the derivative and cofinality calculations be lifted with a modulus depending only on the size of the comparator and the largest marker index it uses?

\end{document}
